\chapter*{习题}
\phantomsection
\addcontentsline{toc}{chapter}{习题}
\setcounter{exercisegroup}{1}
\edef\CurrentExerciseGroup{\arabic{exercisegroup}}
\section*{\texorpdfstring{\S\CurrentExerciseGroup}{第{\CurrentExerciseGroup}节}}
\phantomsection
\addcontentsline{toc}{section}{第{\CurrentExerciseGroup}节习题}

\begin{enumerate}
\def\labelenumi{\arabic{enumi})}
\item
  考察与区间 {[}0, 1{]} 相同的局部紧空间 T、T'，以及 T 与 T' 上分别与 Lebesgue 测度相同的测度 \(\mu, \mu'\)。设 F 是一个 Hilbert 空间，具有可数标准正交基，并将其排成双重序列 \((\mathbf{e}_{mn})\)；设 F' = F 为其对偶。令 \(\mathbf{u}_m(t) = \mathbf{e}_{mn}\)，适用于 \((n-1)2^{-m} \leqslant t < n2^{-m}\) 与 \(1 \leqslant n \leqslant 2^m\)；令 \(\mathbf{f}(t, t') = 2^m \mathbf{u}_m(t)\)，适用于 \(2^{-m} < t' \leqslant 2^{-m+1}\)，最后令 \(\mathbf{f}(t, 0) = 0\) 与 \(\mathbf{f}(1, t') = 0\)。试证 \(\mathbf{f}\) 关于乘积测度 \(\mu \otimes \mu'\) 是标量可积的，但函数 \(t' \mapsto \mathbf{f}(t, t')\) 不是标量 \(\mu'\)-可积的——对任何 \(t \in T\) 都是如此。
\item
  设 \(\mathrm{T}, \mathrm{X}, \mathrm{Y}\) 是三个局部紧空间，其中 \(\mathrm{Y}\) 具有可数基。设 \(\mu\) 是一个测度 \(\geqslant 0\)，定义在 \(\mathrm{T}\) 上；设 \(t \mapsto \lambda_t (t \in \mathrm{T})\) 是一个 \(\mu\)-适宜\(^1\) 的测度族，\(\geqslant 0\)，定义在 \(\mathrm{X}\) 上，并令 \(\nu = \int \lambda_t d\mu(t)\)。设 \(x \mapsto \rho_x (x \in \mathrm{X})\) 是一个 \(\nu\)-适宜的测度族，\(\geqslant 0\)，定义在 \(\mathrm{Y}\) 上。此外，假设还满足下列两个条件之一：a) \(\mathrm{X}\) 是无穷远处可数的；
\end{enumerate}

\begin{enumerate}
\def\labelenumi{\alph{enumi})}
\setcounter{enumi}{1}
\tightlist
\item
  \(\nu\) 是有界的。
\end{enumerate}

在这些条件下，试证存在一个局部 \(\mu\)-可忽略集 \(N'\)，使得对每个 \(t \notin N'\)，族 \(x \mapsto \rho_x\) 都是 \(\lambda_t\)-适宜的；此外，族 \(t \mapsto \int \rho_x d\lambda_t(x)\)（对 \(t \notin N'\) 有定义）是 \(\mu\)-适宜的，且

\[
\int d \mu (t) \int \rho_ {x} d \lambda_ {t} (x) = \int \rho_ {x} d \nu (x).
\]

（为证 \(x \mapsto \rho_x\) 局部几乎处处为标量 \(\lambda_t\)-可积，可利用 §3 第 1 号的引理 1，在情形 \(b\) 中要利用 \(\lambda_t\) 局部几乎处处有界这一事实；用同样的方法可证 \(t \mapsto \int \rho_x d\lambda_t(x)\) 是标量 \(\mu\)-可积的；再利用第 5 号的命题 13。）

\begin{enumerate}
\def\labelenumi{\arabic{enumi})}
\setcounter{enumi}{2}
\tightlist
\item
  设 S、T、Y 是三个无穷远处可数的局部紧空间，且 Y 具有可数基。设 \(\rho\)（相应地 \(\sigma\)）是 S（相应地 T）上的正测度，\(\nu = \rho \otimes \sigma\)，并设 \((\lambda_{s,t})_{(s,t) \in \mathbb{S} \times \mathbb{T}}\) 是 Y 上的一个 \(\nu\)-适宜\(^{1}\)正测度族。那么，存在一个 \(\rho\)-可忽略集 N，使得对每个 \(s \notin N\)，族 \((\lambda_{s,t})_{t \in \mathbb{T}}\) 都是 \(\sigma\)-适宜的；族 \(s \mapsto \int \lambda_{s,t} d\sigma(t)\)（几乎处处有定义）是 \(\rho\)-适宜的，且
\end{enumerate}

\[
\iint \lambda_ {s, t} d \rho (s) d \sigma (t) = \int d \rho (s) \int \lambda_ {s, t} d \sigma (t).
\]

（对空间 \(\mathrm{X} = \mathrm{S} \times \mathrm{T}\) 应用习题 2，并注意 \(\nu = \int (\varepsilon_s \otimes \sigma) d\rho(s)\)。）

\begin{enumerate}
\def\labelenumi{\arabic{enumi})}
\setcounter{enumi}{3}
\tightlist
\item
  设 \(\mathrm{T}, \mathrm{X}, \mathrm{Y}\) 是三个局部紧空间，其中 \(\mathrm{X}\) 是无穷远处可数的。设 \(\mu\) 是 \(\mathrm{T}\) 上的正测度，\((\lambda_t)_{t \in \mathrm{T}}\) 是一个 \(\mu\)-适宜正测度族（在 \(\mathrm{X}\) 上），且 \(\nu = \int \lambda_t d\mu(t)\)。假设下列两个条件之一成立：
\end{enumerate}

\begin{enumerate}
\def\labelenumi{\alph{enumi})}
\item
  Y 具有可数基；
\item
  测度 \(\nu\) 是有界的。
\end{enumerate}

在这些条件下，若 \(\pi\) 是 X 到 Y 中的 \(\nu\)-逆紧映射，则诸 \(t \in \mathbf{T}\) 中使 \(\pi\) 不是 \(\lambda_t\)-逆紧者所成之集是局部 \(\mu\)-可忽略的；此外，Y 上的正测度族 \((\pi(\lambda_t))_{t \in \mathbf{T}}\)（在 T 中局部几乎处处有定义）是 \(\mu\)-适宜的，且

\[
\pi \left(\int \lambda_ {t} d \mu (t)\right) = \int \pi (\lambda_ {t}) d \mu (t).
\]

（在情形 \(a\)，对 \(\rho_x = \varepsilon_{\pi(x)}\) 应用习题 2。在情形 \(b\)，注意 \(\lambda_t\) 局部几乎处处有界，于是可归结为证明映射 \(t \mapsto \pi(\lambda_t)\) 是淡 \(\mu\)-可测的（参见第 V 章 §3 第 2 号命题 4）。给定一个数 \(\varepsilon > 0\)，并设 \(K\) 是 \(T\) 的紧子集，首先注意 \(X\) 中存在递增的紧集序列 \((F_m)\)，使 \(\pi\) 在每个 \(F_m\) 上的限制连续，且令 \(N_m = X - F_m\) 时有 \(\nu(N_m) \leqslant \varepsilon/4^m\)。设 \(A_m\) 是 \(t \in K\) 中使 \(N_m\) 不是 \(\lambda_t\)-可积或使 \(\lambda_t(N_m) > 1/2^m\) 的那些 t 所成的集；若 \(A\) 是诸 \(A_m\) 之并，则 \(\mu(A) \leqslant \varepsilon/2\)。设 \(K_1\) 是 \(K - A\) 的紧子集，使 \(\mu(K - K_1) \leqslant \varepsilon\)，且对每个函数 \(g \in \mathcal{H}(X)\)，\(t \mapsto \int g(x) d\lambda_t(x)\) 在 \(K_1\) 上的限制连续。利用 Urysohn 定理，试证对每个函数 \(f \in \mathcal{H}(Y)\)，在 \(K_1\) 上 \(t \mapsto \int f(\pi(x)) d\lambda_t(x)\) 的限制是连续函数的一致极限。）

*5) 对每个 \(t \in \mathbf{R}\)，令 \(\mathbf{f}(t)\) 为连续函数 \(x \mapsto \exp(-2\pi itx)\)，它是空间 \(\mathcal{C} = \mathcal{C}_{\mathbf{C}}(\mathbf{R})\)（\(\mathbf{R}\) 上的连续复函数空间）中的一个元素。空间 \(\mathcal{C}\) 与空间 \(\mathcal{C}' = \mathcal{C}'_{\mathbf{C}}(\mathbf{R})\)（\(\mathbf{R}\) 上具有紧支集的复测度组成的空间）成对偶；具有紧支集的连续复函数空间 \(\mathcal{K} = \mathcal{K}_{\mathbf{C}}(\mathbf{R})\) 可按典范方式与 \(\mathcal{C}'\) 的一个子空间等同，办法是把函数 \(\varphi \in \mathcal{K}\) 等同于关于 Lebesgue 测度以 \(\varphi\) 为密度的测度。用 \(\mathcal{D}\) 表示由 \(\mathcal{K}\) 中具有紧支集的无穷次可微复函数所构成的子空间。试证：当 \(\mathcal{C}\) 赋予拓扑 \(\sigma(\mathcal{C},\mathcal{C}')\) 或拓扑 \(\sigma(\mathcal{C},\mathcal{K})\) 时，\(\mathbf{f}\) 关于 Lebesgue 测度 \(\mu\) 不是标量可积的；然而对拓扑 \(\sigma(\mathcal{C},\mathcal{D})\)，\(\mathbf{f}\) 是标量 \(\mu\)-可积的，且 \(\int f d\mu\) 是测度 \(\varepsilon_0\)；试证在此情形命题 7 的条件得到验证。*

\begin{enumerate}
\def\labelenumi{\arabic{enumi})}
\setcounter{enumi}{5}
\item
  设 F 是一个可分辨的 Fréchet 空间（TVS，IV，§2，习题 4 b）），F' 是其对偶。试证：若 f 是 T 到 F 中的标量本质可积映射，则 \(\int f d\mu\) 属于 F 的双对偶 F'\,'。（把 F 嵌入 F'\,'；应用定理 1 以及命题 2 与附录的引理。）
\item
  \begin{enumerate}
  \def\labelenumii{\alph{enumii})}
  \tightlist
  \item
    设 \(\mathbf{F} = \overline{\mathcal{K}(\mathbf{N})}\) 是趋于 0 的实数序列组成的 Banach 空间，\(\mathbf{F}' = \mathbf{L}^1 (\mathbf{N})\) 是其对偶（TVS，IV，§1，习题 1）。设 \(\mathbf{f} = (f_n)\) 是 \(\mathrm{I} = [0,1]\) 到 \(\mathbf{F}\) 中的映射，满足 \(f_{n}(t) = n\varphi_{\mathrm{I}_{n}}(t)\)，其中 \(\mathrm{I}_n = [0,1 / n]\)。试证 \(\mathbf{f}\) 关于 Lebesgue 测度 \(\mu\)（在 \(\mathrm{I}\) 上）是标量可积的，但 \(\int \mathbf{f}d\mu\) 是 \(\mathbf{F}''\) 中不属于 \(\mathbf{F}\) 的元素。
  \end{enumerate}
\end{enumerate}

\begin{enumerate}
\def\labelenumi{\alph{enumi})}
\setcounter{enumi}{1}
\item
  设 \(\mathbf{g}\) 是 \(J = [-1, +1]\) 到 \(F\) 中的映射，由条件 \(\mathbf{g}(t) = \mathbf{f}(t)\)（对 \(t \geqslant 0\)）与 \(\mathbf{g}(t) = -\mathbf{f}(-t)\)（对 \(t \leqslant 0\)）定义。则 \(\mathbf{g}\) 关于 \(J\) 上的 Lebesgue 测度是标量可积的，且 \(\int \mathbf{g} d\mu \in F\)，但存在函数 \(h \in \mathcal{L}^{\infty}\) 使 \(\int h\mathbf{g} d\mu \notin F\)。
\item
  由 \(a\) 推出拓扑向量空间 \(F\) 不与任何赋范空间的对偶同构这一事实的一个新证明（TVS，IV，§2，习题 22 c））。
\end{enumerate}

\begin{enumerate}
\def\labelenumi{\arabic{enumi})}
\setcounter{enumi}{7}
\tightlist
\item
  设 \(\mathbf{F}\) 是一个 Hausdorff 局部凸空间，\(\mathbf{f}\) 是 \(\mathbf{T}\) 到 \(\mathbf{F}\) 中的标量局部可积映射，使得对每个紧子集 \(\mathbf{K}\)（\(\mathbf{T}\) 中的）都有 \(\int_{\mathbf{K}} \mathbf{f} d\mu \in \mathbf{F}\)。
\end{enumerate}

\begin{enumerate}
\def\labelenumi{\alph{enumi})}
\item
  试证：若 \(\mathbf{f}\) 是标量本质可积的，且 \(\mathrm{F}\) 是半自反的（或者只需每个对 \(\sigma (\mathrm{F},\mathrm{F}^{\prime})\) 的 Cauchy 序列对该拓扑都在 \(\mathrm{F}\) 中收敛，当 \(\mathrm{T}\) 是无穷远处可数时），则 \(\int g\mathbf{f}d\mu \in \mathrm{F}\) 对每个函数 \(g\in \mathcal{L}^{\infty}\) 成立。
\item
  试证：若 F 是拟完备的，且对 F 上每个连续半范数 q 都有 \(\int^{\bullet} q(\mathbf{f}(t)) d\mu(t) < +\infty\)，则 f 是标量本质可积的，且 \(\int g\mathbf{f} d\mu \in F\) 对每个函数 \(g \in L^{\infty}\) 成立。
\end{enumerate}

（在两种情形中，都考察 T 的紧子集组成的递增有向集。）

\begin{enumerate}
\def\labelenumi{\arabic{enumi})}
\setcounter{enumi}{8}
\item
  设 G、H 是两个 Hausdorff 局部凸空间，U 是 T 到 \(F = \mathcal{L}_{s}(G;H)\) 中的映射。假设 G 是桶式的，H 是拟完备的，U 是 \(\mu\)-可测的，且对 T 的每个紧子集 K，\(U(K)\) 都有界；此外还假设，对 H 上每个连续半范数 q 以及每个 \(x \in G\)，\(\int^{\bullet} q(U(t) \cdot x) d\mu(t) < +\infty\)。在这些条件下，试证 U 是标量本质 \(\mu\)-可积的，且 \(\int U(t) d\mu(t) \in F\)（利用习题 8 b））。
\item
  设 \(G, H\) 是两个 Fréchet 空间，\(G_0\)（相应地 \(H_0\)）是 \(G\)（相应地 \(H\)）的稠密子集，\(t \mapsto \Phi_t\) 是 \(T\) 到空间 \(F = \mathcal{B}(G, H)\)（即 \(G \times H\) 上连续双线性型组成的空间，赋以逐点收敛拓扑）中的映射。假设：\(1^\circ\) 对每个对 \((a, b) \in G_0 \times H_0\)，\(t \mapsto \Phi_t(a, b)\) 都是本质 \(\mu\)-可积的；\(2^\circ\) 若对每对有界子集 \(A \subset G\)、\(B \subset H\)，令 \(q_{A,B}(\Phi) = \sup_{(x, y) \in A \times B} |\Phi(x, y)|\)，其中
\end{enumerate}

对每个 \(\Phi \in \mathcal{B}(\mathrm{G},\mathrm{H})\) 都有 \(\int^{*}q_{\mathrm{A},\mathrm{B}}(\Phi_t)d\mu (t) < + \infty\)。在这些条件下，试证 \(t\mapsto \Phi_t\) 是标量本质 \(\mu\)-可积的，且 \(\int \Phi_t d\mu (t)\in \mathrm{F}\)（利用 Lebesgue 定理，归结为应用定理 1 的推论 1）。\(\mathrm{H} = \mathbf{R}\) 的特殊情形。

\begin{enumerate}
\def\labelenumi{\arabic{enumi})}
\setcounter{enumi}{10}
\item
  设 \(\mu\) 是 \(T = [0,1]\) 上的 Lebesgue 测度，\(F\) 是具有可数标准正交基 \((\mathbf{e}_n)\) 的 Hilbert 空间。设 \(\mathbf{f}\) 是 \(T\) 到 \(F\) 中的映射，满足 \(\mathbf{f}(0) = 0\)，且 \(\mathbf{f}(t) = 2^n \mathbf{e}_n / n\) 在区间 \([2^{-n-1}, 2^{-n}]\) 上成立（\(n \geqslant 0\)）。试证 \(\mathbf{f}\) 是 \(\mu\)-可测的、标量 \(\mu\)-可积的，且 \(\int \mathbf{f} d\mu \in F\)，但 \(\int^* |\mathbf{f}| d\mu = +\infty\)。
\item
  设 \(\mu\) 是 \(T = [0,1]\) 上的 Lebesgue 测度，并设 \(F\) 是一个 Hilbert 空间，具有标准正交基 \((\mathbf{e}_t)_{t \in T}\)，与 \(T\) 等势。
\end{enumerate}

\begin{enumerate}
\def\labelenumi{\alph{enumi})}
\item
  设 \(\mathbf{f}\) 是映射 \(t \mapsto \mathbf{e}_t\)，从 \(\mathrm{T}\) 到 \(\mathrm{F}\) 中。试证 \(\mathbf{f}\) 是标量 \(\mu\)-可忽略的，但不是 \(\mu\)-可测的（当 \(\mathrm{F}\) 赋以弱拓扑 \(\sigma(\mathrm{F}, \mathrm{F}')\) 时；更不用说当 \(\mathrm{F}\) 赋以其原拓扑时了）。
\item
  设 A 是 T 中的一个不可测集（第 IV 章 §4，习题 8）；试证函数 \(\mathbf{g} = \mathbf{f}\varphi_{\mathrm{A}}\) 是标量 \(\mu\)-可忽略的，但数值函数 \(|\mathbf{g}|\) 不是 \(\mu\)-可测的。
\end{enumerate}

\begin{enumerate}
\def\labelenumi{\arabic{enumi})}
\setcounter{enumi}{12}
\item
  设 F 是 Hausdorff 局部凸空间，\(F'\) 是其对偶，q 是 F 上的一个下半连续半范数。试证：若 f 是 T 到 F 中的映射，是 \(\mu\)-可测的（对拓扑 \(\sigma(\mathrm{F},\mathrm{F}')\)），则数值函数 \(q \circ f\) 是 \(\mu\)-可测的。
\item
  设 \(\mu\) 是 T = {[}0,1{]} 上的 Lebesgue 测度；用 F 表示由 T 上有限的 \(\mu\)-可测数值函数组成的 R 上的向量空间，赋以逐点收敛拓扑，这使它成为一个 Hausdorff 局部凸空间。
\end{enumerate}

\begin{enumerate}
\def\labelenumi{\alph{enumi})}
\item
  试证 F 中存在可数稠密子集（考察 T 上连续数值函数组成的 Banach 空间 \(\mathcal{C}(T)\) 中的一个稠密序列）。
\item
  对每个 \(t \in T\)，令 \(\mathbf{f}(t)\) 是 F 的对偶 \(F'\) 中由 \(\langle\mathbf{f}(t), z\rangle = z(t)\)（对一切 \(t \in T\)）定义的元素。试证：当 \(F'\) 赋以拓扑 \(\sigma(\mathrm{F}', \mathrm{F})\) 时，f 是标量 \(\mu\)-可测的，但不是 \(\mu\)-可测的（参见命题 13）。
\end{enumerate}

\begin{enumerate}
\def\labelenumi{\arabic{enumi})}
\setcounter{enumi}{14}
\item
  设 F 是可度量化的局部凸空间，f 是 T 到 F 中的标量 \(\mu\)-可测映射，使得对 T 的每个紧子集 K，存在 F 的可数子集 H，使 \(\mathbf{f}(t) \in \overline{\mathbf{H}}\) 对几乎所有的 \(t \in K\) 成立；试证在这些条件下，f 对 F 的原拓扑是 \(\mu\)-可测的。（在化归到 F 为可分的情形之后，把 F 嵌入可数个赋范空间的乘积中。）
\item
  设 \(\mathbf{F}\) 是 Banach 空间，\(\mathbf{F}'\) 是其对偶。对每个 \(p\)，当 \(1 \leqslant p \leqslant +\infty\) 时，用 \(\Lambda_{\mathbf{F}'}^p(\mathrm{T},\mu)\)（或简记为 \(\Lambda_{\mathbf{F}'}^p\)）表示由这样的映射 \(\mathbf{f}\)——即 \(\mathrm{T}\) 到 \(\mathbf{F}'\) 中的、使得对每个 \(\mathbf{z} \in \mathbf{F}\) 数值函数 \(\langle \mathbf{z},\mathbf{f}\rangle\) 属于 \(\mathcal{L}^p(\mathrm{T},\mu)\) 的映射——组成的集。有 \(\Lambda_{\mathbf{F}'}^p \supset \mathcal{L}_{\mathbf{F}'}^p\)，且 \(\Lambda_{\mathbf{F}'}^1\) 是标量本质 \(\mu\)-可积的、取值于 \(\mathbf{F}'\) 的函数组成的空间（\(\mathbf{F}'\) 赋以 \(\sigma(\mathbf{F}',\mathbf{F})\)）。已知，若用 \(\theta(\mathbf{z})\) 表示 \(\langle \mathbf{z},\mathbf{f}\rangle\) 在 \(\mathrm{L}^p(\mu)\) 中的类，则映射 \(\mathbf{z} \mapsto \theta(\mathbf{z})\) 是连续的（第 4 号，引理 2）；设 \(M_p(\mathbf{f})\) 为其范数；它是空间 \(\Lambda_{\mathbf{F}'}^p\) 上的一个半范数。为使 \(M_p(\mathbf{f}) = 0\)，必要且充分的是 \(\mathbf{f}\) 为标量局部可忽略的。
\end{enumerate}

\begin{enumerate}
\def\labelenumi{\alph{enumi})}
\item
  为使 \(\mathbf{f} \in \Lambda_{\mathbf{F}'}^p\)，必要且充分的是：对每个数值函数 \(g \in \mathcal{L}^q\)，函数 \(g\mathbf{f}\) 都是标量本质 \(\mu\)-可积的；此时有 \(\mathrm{M}_1(g\mathbf{f}) \leqslant \mathrm{M}_p(\mathbf{f})\mathrm{N}_q(g)\)。
\item
  在空间 \(\Lambda_{\mathbf{F}'}^1\) 中，试证半范数 \(\mathbf{M}_1\) 等价于半范数 \(\mathrm{M}_1'(\mathbf{f}) = \sup |\int_{\mathbf{A}} \mathbf{f} d\mu|\)，其中 A 取遍 T 的可测子集所成的集；更精确地，\(\mathrm{M}_1' \leqslant \mathrm{M}_1 \leqslant 2\mathrm{M}_1'\)。由此推出 \(\mathrm{M}_p\) 是等价于 \(\mathrm{M}_p'(\mathbf{f}) = \sup_{\mathrm{N}_q(g) \leqslant 1} \mathrm{M}_1'(g\mathbf{f})\) 的半范数。
\item
  取 F 为一个具有可数标准正交基的 Hilbert 空间，取 \(\mu\) 为 \(T = [0,1]\) 上的 Lebesgue 测度。试证存在 T 到 F 中的一个映射，它不可积但标量可积，并且是可测阶梯函数序列在半范数 \(M_1\) 下的极限（参见习题 11）；由此推出：在空间 \(\mathcal{L}_{\mathrm{F}}^1\)（即 \(\mu\)-可积函数组成的空间）上，由半范数 \(M_1\) 定义的拓扑严格粗于由 \(N_1\) 定义的拓扑。对半范数 \(M_p\) 与 \(N sub p\)（\(1 \leqslant p < +\infty\)）有类似的结果。
\end{enumerate}

\(d\)）试证：在空间 \(\mathcal{L}_{\mathrm{F}^{\prime}}^{\infty}\) 上，\(\mathrm{M}_{\infty} = \mathrm{N}_{\infty}\)

\begin{enumerate}
\def\labelenumi{\alph{enumi})}
\setcounter{enumi}{4}
\tightlist
\item
  取 F 为一个具有可数标准正交基的 Hilbert 空间，该基排成双重序列 \((\mathbf{e}_{mn})\)。设 \(\mu\) 是 \(T = [0,1]\) 上的 Lebesgue 测度。设 \(\mathbf{u}_m\) 是 T 到 F 中的映射，满足 \(\mathbf{u}_m(t) = \mathbf{e}_{mn}\)（对 \((n - 1)2^{-m} \leqslant t < n2^{-m}\) 及 \(1 \leqslant n \leqslant 2^m\)），且 \(\mathbf{u}_m(1) = 0\)；令 \(\mathbf{f}_n = \sum_{i=0}^{n} \mathbf{u}_i\)。试证：对 \(1 \leqslant p < +\infty\)，\((\mathbf{f}_n)\) 是 \(\Lambda_{\mathrm{F}'}^p\) 中的 Cauchy 序列，但不收敛于 \(\Lambda_{\mathrm{F}'}^p\) 中的任何函数。
\end{enumerate}

\begin{enumerate}
\def\labelenumi{\alph{enumi})}
\setcounter{enumi}{5}
\item
  设 \((\mathbf{f}_{n})\) 是空间 \(\Lambda_{F'}^{p}\)（\(1 \leqslant p \leqslant +\infty\)）中的 Cauchy 序列；假设存在 T 到 \(F'\) 中的映射 f，使得对每个 \(z \in F\)，函数序列 \(\langle z, f_{n} \rangle\) 都依测度收敛于 \(\langle z, f \rangle\)（第 IV 章 §5 第 11 号）。试证 \(f \in \Lambda_{F'}^{p}\)，且序列 \((\mathbf{f}_{n})\) 在 \(\Lambda_{F'}^{p}\) 中以 f 为极限。
\item
  取 F 为绝对收敛级数组成的 Banach 空间 \(\mathrm{L}^{1}(\mathbf{N})\)；其对偶为 \(\mathbf{F}^{\prime} = \mathrm{L}^{\infty}(\mathbf{N})\)（TVS，IV，§1，习题 1）。设 \(\mu\) 是 T = {[}0,1{]} 上的 Lebesgue 测度。设 f 是 T 到 \(F^{\prime}\) 中的映射，满足 \(\mathbf{f}(0) = 0\)，且对 \(2^{-n-1} < t \leqslant 2^{-n}\)，\(\mathbf{f}(t)\) 是这样的序列：它的各项皆为零，除了指标为 n 的项等于 \(2^{n/p}\)（\(1 \leqslant p < +\infty\)）。试证 f 对 \(F^{\prime}\) 上的强拓扑是可测的，且属于 \(\Lambda_{\mathrm{F}'}^p\)，但任何阶梯函数序列 \((\mathbf{f}_n)\) 都不趋于 \(\mathbf{f}\)（在空间 \(\Lambda_{\mathrm{F}'}^p\) 中）。
\end{enumerate}

¶ 17) 设 \(\mu\) 是 \(T = [0,1]\) 上的 Lebesgue 测度，\(F\) 是 Banach 空间 \(L^1(\mathbf{N})\)，\(F' = L^\infty(\mathbf{N})\) 是其对偶。对每个 \(t \in [0,1[\)，设 \(t = \sum_{n=0}^\infty \xi_n 2^{-n}\) 为 \(t\) 的二进展开；用 \(f(t)\) 表示序列 \((\xi_n) \in F'\)，并令 \(f(1) = 0\)。

\begin{enumerate}
\def\labelenumi{\alph{enumi})}
\item
  试证：对 \(1 \leqslant p \leqslant +\infty\)，函数 \(\mathbf{f}\) 属于空间 \(\Lambda_{\mathrm{F}'}^p\)（参见习题 16）。
\item
  试证 \(\mathbf{f}\) 对拓扑 \(\sigma(\mathrm{F}',\mathrm{F})\) 是可测的，且对该拓扑存在一个阶梯函数序列几乎处处收敛于 \(\mathbf{f}\)。
\item
  试证 \(\mathbf{f}\) 对 \(\mathbf{F}'\) 上的强拓扑不是可测的（注意 \(|\mathbf{f}(t) - \mathbf{f}(t')| = 1\)，只要 \(0 \leqslant t < t' < 1\)）。
\item
  试证：在空间 \(\Lambda_{\mathrm{F}'}^p\)（\(1 \leqslant p < +\infty\)）中，\(\mathbf{f}\) 不是任何阶梯函数序列的极限。（可归结为只考虑这样的阶梯函数：它们是区间的特征函数的线性组合（系数取自 \(\mathbf{F}'\)），其中区间的端点形如 \(k/2^n\)；若 \(\mathbf{g}\) 是这样一个阶梯函数，试证在习题 16 b) 的记号下有 \(\mathrm{M}_1'(\mathbf{f} - \mathbf{g}) \geqslant 1/4\)。）
\item
  试证不存在任何映射 \(\mathbf{h}\)——\(\mathrm{T}\) 到 \(\mathrm{F}'\) 中的、对强拓扑可测的映射——使 \(\mathbf{f} - \mathbf{h}\) 是标量可忽略的。（注意对这样的映射，必然有 \(\mathrm{N}_{\infty}(\mathbf{h}) \leqslant 1\)；然后由此得出与 \(d\) 的结果矛盾。）
\end{enumerate}

\begin{enumerate}
\def\labelenumi{\arabic{enumi})}
\setcounter{enumi}{17}
\tightlist
\item
  试证：当把 \(\mathbf{f}\) 为 \(\mu\)-可测（对 \(\mathbf{F}\) 的原拓扑）这一条件换成 \(\mathbf{f}\) 为 \(\mu\)-可测（对弱化拓扑 \(\sigma(\mathrm{F},\mathrm{F}')\)）这一条件时，命题 8 仍然成立。（利用第 IV 章 §4 习题 19 c）。）
\end{enumerate}

¶ 19) 设 f 是 T 到 F 中的标量本质 μ-可积映射。对每个 \(z' \in F'\)，用 \(u_{\mathbf{f}}(\mathbf{z}')\) 表示 \(L^{1}(\mu)\) 中由函数 \(\langle f, z' \rangle\) 所确定的类。若对每个函数 \(g \in L^{\infty}(\mu)\) 都有 \(\int g f d\mu \in F\)，则称 f 是标量良可积的。

\begin{enumerate}
\def\labelenumi{\alph{enumi})}
\item
  为使 \(\mathbf{f}\) 是标量良可积的，必要且充分的是 \(u_{\mathbf{f}}\) 对拓扑 \(\sigma(\mathrm{F}',\mathrm{F})\) 与 \(\sigma(\mathrm{L}^1,\mathrm{L}^\infty)\) 连续。当此条件成立时，在 \(u_{\mathbf{f}}\) 之下，\(\mathrm{F}'\) 的每个等度连续子集的像是 \(\mathrm{L}^1\) 中的相对弱紧子集。特别地：\(\alpha\) 对每个等度连续子集 \(\mathrm{H}'\)（\(\mathrm{F}'\) 的）以及每个 \(\varepsilon >0\)，存在紧子集 \(\mathrm{L}\)（在 \(\mathrm{T}\) 中），使 \(\int_{\mathrm{T}-\mathrm{L}}|\langle\mathbf{f},\mathbf{z}'\rangle|d\mu\leqslant\varepsilon\) 对一切 \(\mathbf{z}'\in\mathrm{H}'\) 成立；\(\beta\) 对每个等度连续子集 \(\mathrm{H}'\)（\(\mathrm{F}'\) 的）以及每个 \(\varepsilon>0\)，存在 \(\eta>0\)，使得对每个开集 \(\mathrm{U}\subset\mathrm{T}\)（测度 \(\mu(\mathrm{U})\leqslant\eta\)），都有 \(\int_{\mathrm{U}}|\langle\mathbf{f},\mathbf{z}'\rangle|d\mu\leqslant\varepsilon\) 对一切 \(\mathbf{z}'\in\mathrm{H}'\) 成立（第 V 章 §5，习题 15）。
\item
  假设条件 \(\alpha\) 与 \(\beta\)（\(a\) 中的）都成立，此外还存在子空间 \(\mathrm{H} \subset \mathrm{L}^{\infty}\)，它对拓扑 \(\sigma(\mathrm{L}^{\infty}, \mathrm{L}^{1})\) 稠密，使得对每个函数 \(g \in \mathcal{L}^{\infty}\)，只要其类属于 \(\mathrm{H}\)，就有 \(\int g\mathbf{f}d\mu \in \mathrm{F}\)。试证在这些条件下，若再假设 \(\mathrm{F}\) 是拟完备的，则 \(\mathbf{f}\) 是标量良可积的。（先化归到 \(\mathrm{F}\)（对其原拓扑）完备的情形，并注意 \(\dot{g} \mapsto \int g\mathbf{f}d\mu\) 是 \(\mathrm{L}^{\infty}\) 到 \(\mathrm{F}'*\) 中的连续线性映射，其中 \(\mathrm{L}^{\infty}\) 赋以拓扑 \(\tau(\mathrm{L}^{\infty}, \mathrm{L}^{1})\)，\(\mathrm{F}'*\) 赋以在 \(\mathrm{F}'\) 的等度连续子集上一致收敛的拓扑。）
\item
  设 \(\mathbf{f}\) 是标量良可积的。试证对每个 Cauchy 序列 \((\mathbf{z}_n')\)（对拓扑 \(\sigma(\mathrm{F}',\mathrm{F})\)），序列 \((u_{\mathbf{f}}(\mathbf{z}_n'))\) 在 \(\mathrm{L}^1\) 中强收敛（利用第 IV 章 §5 习题 25 c) 与第 V 章 §5 习题 16 b)）。由此推出：若 A 是 \(\mathrm{L}^{\infty}\) 的单位球在映射 \(\dot{g} \mapsto \int g\mathbf{f}d\mu\) 下的像，则每个 Cauchy 序列 \((\mathbf{z}_n')\)（对拓扑 \(\sigma(\mathrm{F}',\mathrm{F})\)）都是对在 A 上一致收敛的拓扑的 Cauchy 序列。
\end{enumerate}

\(d)\) 由 \(c)\) 推出：若 \(\mathbf{f}\) 是标量良可积的，则 \(u_{\mathbf{f}}\) 把 \(\mathrm{F}'\) 的每个对 \(\sigma (\mathrm{F}',\mathrm{F})\) 等度连续且可度量化的子集映为 \(L^{1}\) 中的相对强紧集。特别地，若 F 是 Hausdorff 局部凸空间 G 的对偶 \(G'\)，F 赋以一个与 F 和 G 之间的对偶相容的拓扑，且 \(G'\) 中存在可数的强稠密集，则 \(u_{f}\) 把 \(F' = G\) 的每个有界子集映为 \(L^{1}\) 中相对强紧的集（把 G 嵌入其双对偶）。

\begin{enumerate}
\def\labelenumi{\arabic{enumi})}
\setcounter{enumi}{19}
\tightlist
\item
  设 F 是 Hausdorff 局部凸空间，\(F'\) 是其对偶，\(S'\) 是由这样的子集 \(A' \subset F'\) 组成的集：\(A'\) 的每个点序列都有对 \(\sigma(\mathrm{F}',\mathrm{F})\) 收敛的子序列。试证：对 F 的有界子集 A，下列条件等价：\(\alpha\) A 对 \(S'\)-拓扑是预紧的；\(\beta\) 每个对 \(\sigma(\mathrm{F}',\mathrm{F})\) 收敛的序列都在 A 上一致收敛。（设 S 是由 F 与 A 的有限子集组成的集；注意 \(\beta\) 蕴含每个集 \(A' \in S'\) 对 S-拓扑是预紧的，这里要用 GT，II，§4 习题 6。然后应用 TVS，IV，§1 习题 4。）
\end{enumerate}

¶ 21) 设 f 是 T 到拟完备 Hausdorff 局部凸空间 F 中的标量良可积映射。

\begin{enumerate}
\def\labelenumi{\alph{enumi})}
\item
  试证在下列每种情形中，映射 \(\dot{g} \mapsto \int g\mathbf{f}  d\mu\) 把 \(\mathrm{L}^{\infty}\) 的单位球映为 \(\mathbf{F}\) 中对原拓扑的相对紧子集：\(1^{\circ}\) \(\mathbf{F}\) 中存在可数稠密集；\(2^{\circ}\) \(\mathbf{F}\) 是一个对偶 \(\mathbf{G}'\)——其中 \(\mathbf{G}\) 或者可度量化，或者是可度量化空间序列的严格直接极限——且 \(\mathbf{F}\) 赋以拓扑 \(\tau (\mathbf{G}',\mathbf{G})\)。（利用习题 19 c) 与习题 20；在情形 \(2^{\circ}\)，用 Šmulian 定理（TVS，IV，§5，习题 2 c））。）
\item
  试证 a) 的结论在下述情况下仍然成立：不对 F 作新的假设，而假设 f 是可测的。（把 F 嵌入一族 Banach 空间的乘积，化归到 F 是 Banach 空间的情形；然后利用习题 19 a)，化归到 T 是紧的情形；再应用 a) 的情形 \(1^{\circ}\)。
\end{enumerate}

\begin{enumerate}
\def\labelenumi{\arabic{enumi})}
\setcounter{enumi}{21}
\tightlist
\item
  在 Hausdorff 局部凸空间 \(\mathbf{F}\) 中，设 \((\mathbf{z}_{\alpha})_{\alpha \in \mathbf{A}}\) 是这样的族：映射 \(\alpha \mapsto \mathbf{z}_{\alpha}\) 当 \(\mathbf{A}\) 赋以由每点质量 \(+1\) 定义的测度时是标量良可积的（习题 19）。
\end{enumerate}

\begin{enumerate}
\def\labelenumi{\alph{enumi})}
\item
  试证族 \((\mathbf{z}_{\alpha})\) 对每个与 F 和 \(F'\) 之间的对偶相容的拓扑 T 都是可和的（利用习题 19 a)）。当 F 对 T 拟完备时，逆命题成立。
\item
  设 \(F = G'\)，\(F' = G\)，其中 G 是这样的 Hausdorff 局部凸空间：\(G'\) 中存在可数的强稠密集。试证 \((\mathbf{z}_{\alpha})\) 对 \(G'\) 上的强拓扑是可和的（该拓扑不一定与 \(G'\) 与 G 之间的对偶相容（参见习题 19d））。试证这一结果不能推广到 \(G = L^{1}(\mathbf{N})\)、\(G' = L^{\infty}(\mathbf{N})\) 的情形。
\end{enumerate}

\begin{enumerate}
\def\labelenumi{\arabic{enumi})}
\setcounter{enumi}{22}
\tightlist
\item
  设 \(\mu\) 是 \(T\) 上的正测度，由可数个紧集的并所承载。
\end{enumerate}

\begin{enumerate}
\def\labelenumi{\alph{enumi})}
\tightlist
\item
  设 H 是由 \(\mu\)-可测数值函数组成的集，使得对每个 \(t \in \mathbf{T}\)，\(\sup_{f \in \mathrm{H}} f(t) < +\infty\)。试证存在一个有限值的 \(\mu\)-可测函数 \(g\)，使得
\end{enumerate}

\(f(t) \leqslant g(t)\) 对每个函数 \(f \in \mathrm{H}\) 都几乎处处成立。（化归到 \(\mu\) 有界且 \(\mathbf{H}\) 中函数取值于 \([-1, +1]\) 的情形，这可借助 \(\overline{\mathbf{R}}\) 到该区间上的一个递增同胚做到。然后考虑 \(\widetilde{g}\)，即函数的类所成之集在 \(\mathrm{L}^1(\mu)\) 中的上确界，其中函数取自 \(\mathbf{H}\)（第 IV 章 §3，第 6 号，命题 14），并注意存在 \(\mathbf{H}\) 中的函数序列几乎处处收敛于 \(g\)。）

\begin{enumerate}
\def\labelenumi{\alph{enumi})}
\setcounter{enumi}{1}
\item
  设 \(\mathbf{f}\) 是一个标量 \(\mu\)-可测映射，定义在 \(\mathrm{T}\) 上、取值于 Hausdorff 局部凸空间 \(\mathrm{F}\)。对每个子集 \(\mathrm{B}'\)（\(\mathrm{F}'\) 的、对 \(\sigma(\mathrm{F}',\mathrm{F})\) 有界的），试证存在可测的有限值数值函数 \(g_{\mathrm{B}'}\)，使得对每个 \(\mathbf{z}' \in \mathrm{B}'\)，\(|\langle \mathbf{f}(t),\mathbf{z}' \rangle| \leqslant g_{\mathrm{B}'}(t)\) 几乎处处成立。
\item
  设 \(\mathbf{F}\) 是 Fréchet 空间，\(\mathbf{f}\) 是一个标量 \(\mu\)-可测映射，定义在 \(\mathbf{T}\) 上、取值于 \(\mathbf{F}\)。试证：对每个紧子集 \(\mathbf{K}\)（\(\mathbf{T}\) 的）以及每个 \(\varepsilon > 0\)，存在紧集 \(\mathrm{K}' \subset \mathrm{K}\)，使 \(\mu(\mathrm{K} - \mathrm{K}') \leqslant \varepsilon\) 且 \(\mathbf{f}\) 在 \(\mathrm{K}'\) 上的限制有界（利用 \(b\)））。
\end{enumerate}

\begin{enumerate}
\def\labelenumi{\arabic{enumi})}
\setcounter{enumi}{23}
\item
  设 F 是完备的 Hausdorff 局部凸空间，包含可数稠密集。设 f 是 T 到 F 中的映射，标量 \(\mu\)-可测且标量本性有界；试证：对每个函数 \(g \in L^{1}\)，gf 是标量可积的且 \(\int g f d\mu \in F\)。（利用 TVS，III，§3，第 6 号定理 2 的推论 1，注意 \(F'\) 的等度连续子集对 \(\sigma(F', F)\) 可度量化，并应用 Lebesgue 定理。）
\item
  设 F 是可分 Fréchet 空间，其中每个对 \(\sigma(F, F')\) 的 Cauchy 序列都对这一拓扑收敛（例如空间 \(L^1(T', \nu)\)，其中 \(T'\) 具有可数基（第 V 章 §5，习题 16 b）））。试证 T 到 F 中的每个标量本质可积映射 f 都是标量良可积的。（先证明对每个具有紧支集的函数 \(g \in \mathcal{L}^\infty(\mu)\)，有 \(\int g\mathbf{f} d\mu \in F\)，这要注意到 f 是 \(\mu\)-可测的（第 5 号，命题 12），并应用第 2 号的命题 8 与假设。然后借助假设由此推出：\(\int_A g\mathbf{f} d\mu \in F\) 对每个函数 \(g \in \mathcal{L}^\infty\) 以及每个为可数个紧集之并的集 A 都成立。最后，用习题 19 的记号，证明在 \(u_{\mathbf{f}}\) 之下，\(F'\) 的每个等度连续子集的像都在 \(L^1\) 中相对弱紧，这里要用 Eberlein 定理（TVS，IV，§5，第 3 号，定理 1）以及 \(F'\) 的等度连续子集对 \(\sigma(F', F)\) 可度量化这一事实。）
\item
  设 \((f_n)\) 是 \(\mu\)-可积函数序列，定义在 \(\mathbf{T}\) 上，使得对几乎所有 \(t \in \mathbf{T}\)，\(\sup_{n} |f_n(t)| < +\infty\) 且 \(\sup_{n} \int |f_n(t)| d\mu(t) = +\infty\)。试证存在标量序列 \((c_n)\)，使 \(\sum_{n} |c_n| < +\infty\) 且 \(\sum_{n} c_n f_n\) 不是 \(\mu\)-可积的（把 Gelfand-Dunford 定理应用于空间 \(\mathbf{L}^1(\mathbf{N}) = \mathbf{F}\)。）
\item
  设 E 是可分实 Banach 空间，\(G = E'\) 是其对偶，\(G'\) 是 E 的子空间，在 E 中稠密、异于 E 且是桶式的（TVS，III，§4，习题 4）。试证：若 G 赋以拓扑 \(\sigma(G, G')\)，其对偶 \(G' = \mathcal{L}(G; \mathbf{R})\) 赋以紧收敛拓扑，则 \(G'\) 不是拟完备的，且存在序列 \((x_n)\) 在 \(G'\) 中趋于 0，以及一个数序列 \((\lambda_n)\)——这些数 \textgreater0 且和有限——使得通项为 \(\lambda_n x_n\) 的级数在 \(G'\) 中不收敛。（先注意 G 的每个对 \(\sigma(G, G')\) 有界的子集对 G 上的范数有界。取一点 \(a \in E\) 不属于 \(G'\)，再取点序列 \((a_n)\)（\(G'\) 中的点），它在 E 中收敛于 a 且满足 \(\|a_{n+1} - a_n\| \leqslant 1/2^n\)。）
\end{enumerate}

\setcounter{exercisegroup}{2}
\edef\CurrentExerciseGroup{\arabic{exercisegroup}}
\section*{\texorpdfstring{\S\CurrentExerciseGroup}{第{\CurrentExerciseGroup}节}}
\phantomsection
\addcontentsline{toc}{section}{第{\CurrentExerciseGroup}节习题}

\begin{enumerate}
\def\labelenumi{\arabic{enumi})}
\item
  试证：为使数值函数 \(f\)（\(T\) 上的）对 \(T\) 上每个正测度都本质可积，必要且充分的是它有界、有紧支集，且对 \(T\) 上每个正测度都可测。
\item
  设 \(G\) 是桶式空间，\(H\) 是半自反空间，\(F\) 是空间 \(\mathcal{L}_s(G;H)\)。设 \(m\) 是 \(T\) 上的向量测度，取值于 \(F\)。试证：对每个数值函数 \(f\)（关于 \(m\) 本质可积的），有 \(\int fd\mathbf{m} \in F\)。（注意 \(F\) 是半自反的。）\(H = R\) 的情形。
\item
  设 \(\mathbf{m}\) 是 \(\mathrm{T}\) 上的向量测度，取值于 \(\mathrm{F}\)。试证：对每个数值函数 \(f\)（关于 \(\mathbf{m}\) 本质可积的），在下列两种情形的每一种中都有 \(\int f dm \in \mathrm{F}''\)：\(a)\) \(\mathrm{F}\) 是可分辨的（TVS，IV，§2，习题 4）；\(b)\) \(\mathrm{F}\) 是 Fréchet 空间且 \(\mathrm{T}\) 无穷远处可数。（在两种情形中，都把 \(\mathrm{F}\) 嵌入 \(\mathrm{F}''\) 并应用命题 3 的推论 1 的方法；在情形 \(b\) ，用 TVS，IV，§3，第 3 号命题 3 的推论。）
\item
  设 \(\mathbf{m}\) 是 \(\mathrm{T}\) 上的向量测度，取值于 \(\mathrm{F}\)。试证：对每个数值函数 \(f \geqslant 0\)（下半连续且关于 \(\mathbf{m}\) 本质可积的），有 \(\int f d\mathbf{m} \in \mathrm{F}''\)（利用第 IV 章 §1，第 1 号定理 1）。
\item
  设 F 是半自反局部凸空间，m 是 T 上取值于 F 的测度。试证：对每个关于 m 局部可积的数值函数 g，映射 \(f \mapsto \int fg dm\)（\(\mathcal{K}(T)\) 到 F 中的映射）是 T 上的测度，记作 \(g \cdot m\)；为使数值函数 h 关于 \(g \cdot m\) 本质可积，必要且充分的是 gh 关于 m 本质可积，此时有 \(\int h d(g \cdot m) = \int hg dm\)。
\item
  设 F 是半自反局部凸空间，带有一个 R 上的代数结构，使得映射 \((u,v)\mapsto uv\)（\(F\times F\) 到 F 中的映射）对每个变量分别连续。设 m 是 T 上取值于 F 的测度，使得 \(\mathbf{m}(fg)=\mathbf{m}(f)\mathbf{m}(g)\) 对 \(\mathcal{K}(\mathrm{T})\) 中的 f 和 g 成立。试证：若数值函数 f,g 使 f、g 与 fg 都属于 \(\mathcal{L}(\mathbf{m})\)，则 \(\int fg dm=(\int fd\mathbf{m})(\int gd\mathbf{m})\)。（先处理 \(g\in\mathcal{K}(\mathrm{T})\) 的情形，考虑两个测度 \(f\mapsto\mathbf{m}(fg)\) 与 \(f\mapsto\mathbf{m}(f)\mathbf{m}(g)\)；然后类似地处理一般情形并利用习题 5。）\(F=\mathcal{L}(\mathrm{G})\) 的情形，其中 G 自反，F 赋以逐点收敛拓扑或弱逐点收敛拓扑。
\item
  设 \(\mu\) 是 T 上的正测度，m 是 T 上取值于 F 的测度，以 \(\mu\) 为基、以 f 为关于 \(\mu\) 的密度。设 \(q\) 是 F 上的下半连续半范数。试证：若对 T 的每个紧子集 K，\(\int_{K}^{*}(q \circ f) d\mu\) 有限，则对每个 \(\mu\)-可测函数 \(h \geqslant 0\)，不等式 \(\int^{\bullet} h dq(m) \leqslant \int^{\bullet} (q \circ f) h d\mu\) 成立。（利用第 V 章（第 1 版）§2，第 2 号引理 1。）²
\item
  设 \(\mathbf{m}\) 是 \(\mathbf{T}\) 上取值于 \(\mathbf{F}\) 的向量测度，是 \(q\)-可优超的（对一个下半连续半范数 \(q\)，在 \(\mathbf{F}\) 上）。试证：对每个函数 \(f \geqslant 0\)（属于 \(\mathcal{H}(\mathbf{T})\)），
\end{enumerate}

\[
\int f d (q \circ \mathbf {m}) = \sup \sum_ {i} q (\mathbf {m} (f _ {i})),
\]

其中 \((f_i)\) 取遍 \(\mathcal{H}(\mathrm{T})\) 中元素的所有满足 \(\sum_{i}|f_i| \leqslant f\) 的有限序列。（按第 II 章 §2，第 2 号定理 1 的方式论证。）

\begin{enumerate}
\def\labelenumi{\arabic{enumi})}
\setcounter{enumi}{8}
\tightlist
\item
  设 T 是无穷远处可数的局部紧空间，F 是包含可数稠密集的 Hausdorff 桶式空间，m 是 T 上取值于 F 的弱对偶 F' 的向量测度。试证：存在 T 上的正测度 \(\mu\)，使 m 标量意义以 \(\mu\) 为基。（利用第 V 章 §5，第 6 号命题 12，以及 Lebesgue--Nikodym 定理（第 V 章 §5，第 5 号定理 2）推论 5 的条件 5)。）
\end{enumerate}

¶ 10) 设 K 是紧空间。对 K 的每个 Hausdorff 商空间 \(K_{1}\)，把 Banach 空间 \(\mathcal{C}(K_{1})\) 等同于 \(\mathcal{C}(K)\) 的一个闭子空间，借助的是单射 \(f \mapsto f \circ \pi\)，其中 \(\pi\) 是 K 到 \(K_{1}\) 上的典范映射。

\begin{enumerate}
\def\labelenumi{\alph{enumi})}
\item
  设 \(u\) 是 \(\mathcal{C}(\mathrm{K})\) 到拟完备 Hausdorff 局部凸空间 \(\mathrm{F}\) 中的连续线性映射。为使 \(u\) 把 \(\mathcal{C}(\mathrm{K})\) 的单位球变为 \(\mathrm{F}\) 中对 \(\sigma(\mathrm{F},\mathrm{F}')\) 相对紧的子集，只需：对每个可度量化商空间 \(\mathrm{K}_1\)（\(\mathrm{K}\) 的），\(u\) 在 \(\mathcal{C}(\mathrm{K}_1)\) 上的限制具有该性质。（利用 Eberlein 定理（TVS，IV，§5，第 3 号，定理 1），并注意每个序列 \((f_n)\)（在 \(\mathcal{C}(\mathrm{K})\) 中）都含于某个子空间 \(\mathcal{C}(\mathrm{K}_1)\)，其中 \(\mathrm{K}_1\) 是 \(\mathrm{K}\) 的可度量化商空间；为此，考虑连续映射 \(x \mapsto (f_n(x))\)（\(\mathrm{K}\) 到 \(\mathbf{R}^{\mathbf{N}}\) 中的）。）
\item
  对 K 的每个 Hausdorff 商空间 \(K_{1}\)，\(\mathcal{C}(K_{1})\) 到 \(\mathcal{C}(K)\) 中的典范单射的转置是映射 \(\mu\mapsto\pi(\mu)\)（\(\mathcal{M}(K)\) 到 \(\mathcal{M}(K_{1})\) 中的映射）。用 \((\mathcal{M}(\mathrm{K}))'\) 表示 Banach 空间 \(\mathcal{M}(\mathrm{K})\)（即 \(\mathcal{C}(\mathrm{K})\) 的双对偶）的对偶。为使 \(\mathcal{M}(\mathrm{K})\) 的子集 A 对 \(\sigma(\mathcal{M}(\mathrm{K}),(\mathcal{M}(\mathrm{K}))')\) 相对紧，必要且充分的是：对 K 的每个可度量化商空间 \(K_{1}\)，A 在 \(\mathcal{M}(\mathrm{K}_{1})\) 中的典范像对 \(\sigma(\mathcal{M}(\mathrm{K}_{1}),(\mathcal{M}(\mathrm{K}_{1}))')\) 相对紧。（注意 A 是强有界的；于是可设 A 是凸的、均衡的且对淡拓扑 \(\sigma(\mathcal{M}(\mathrm{K}),\mathcal{C}(\mathrm{K}))\) 闭（在 \(\mathcal{M}(\mathrm{K}_{1})\) 中应用第 IV 章 §7 习题 10）；因此 A 是对偶 \(F'\)（某个 Banach 空间 F 的对偶）的单位球在一个连续线性映射 u（\(\mathcal{C}(\mathrm{K})\) 到 F 中的）的转置下的像（考虑 A 在 \(\mathcal{C}(\mathrm{K})\) 中的极）；然后应用 a）。）
\end{enumerate}

¶ 11) 设 F、G 是两个 Hausdorff 局部凸空间，其中 G 假设为拟完备的；设 u 是 F 到 G 中的连续线性映射；其转置 \(^{t}u\) 是强连续的，故转置 \(^{t}(^{t}u)\) 是 F'\,' 到 G'\,' 中的强连续且弱连续的线性映射。试证下列条件等价：

\(\alpha)\) \(u\) 把 \(\mathbf{F}\) 的每个有界子集变为 \(\mathbf{G}\) 中对 \(\sigma(\mathbf{G},\mathbf{G}')\) 相对紧的子集。

\(\beta)^{t}(^{t}u)\) 把 \(\mathbf{F}^{\prime \prime}\) 映到 G 中。

\(\gamma)\ ^{t}u\) 对拓扑 \(\sigma(\mathrm{G}^{\prime},\mathrm{G})\) 与 \(\sigma(\mathrm{F}^{\prime},\mathrm{F}^{\prime\prime})\) 连续。

\(\delta)^{t}u\) 把 \(\mathrm{G}'\) 的每个等度连续子集变为对 \(\sigma (\mathbf{F}',\mathbf{F}'')\) 相对紧的集。

（利用下述事实：F 的每个有界子集在 \(F''\) 中对 \(\sigma(F'', F')\) 相对紧。为看出 \(\delta\) 蕴含 \(\gamma\)，先考虑 G 完备的情形，并用 TVS，III，§3，第 6 号定理 2 的推论 1。为过渡到一般情形，把 G 嵌入其完备化空间，并注意 \(F''\) 的每个点都在 F 的某个有界子集的闭包中（对拓扑 \(\sigma(F'', F')\)）。）

¶ 12) a) 设 K 是紧空间，E = E(K) 是 R \(^{K}\) 的由开集特征函数的线性组合所构成的子空间；把 E 等同于 C(K) 的双对偶的一个子空间。试证：对 σ(M(K),E) 的 Cauchy 序列对 σ(M(K),(M(K))') 收敛（利用第 V 章 §5 的命题 12 与习题 15）。由此推出：M(K) 的每个对拓扑 σ(M(K),E) 相对紧的子集 A，都对拓扑 σ(M(K),(M(K))') 相对紧。（利用习题 10 b)，化归到 K 可度量化的情形。应用第 V 章 §5 习题 13 证明 A 有界。注意在 A 上，淡拓扑 σ(M(K),C(K)) 与拓扑 σ(M(K), \(\overline{E}\)) 相同，其中 \(\overline{E} \supset C(K)\) 是 E 在 M(K) 的强对偶中的闭包；由此并利用 C(K) 可分这一事实，推出 A 的点的每个序列都有对 σ(M(K),E) 的 Cauchy 子序列；最后援引 Eberlein 定理完成证明。）

\begin{enumerate}
\def\labelenumi{\alph{enumi})}
\setcounter{enumi}{1}
\item
  设 T 是局部紧空间，m 是 T 上取值于拟完备 Hausdorff 局部凸空间 F 的向量测度。假设对 T 的每个紧子集 K，\(\int_{K} dm \in F\)；试证：对 T 上每个有紧支集的有界 Borel 函数 f，有 \(\int f dm \in F\)，且在映射 \(f \mapsto \int f dm\) 之下，支集含于 T 的紧子集 K 且范数 \(\|f\| \leqslant 1\) 的 Borel 函数所成之集的像对 \(\sigma(F, F')\) 相对紧（利用 a) 与习题 11，把它应用于 \(u : f \mapsto \int f dm\)，其中 f 取遍 T 的紧子集的特征函数的线性组合所成之集）。
\item
  设 F 是拟完备的 Hausdorff 局部凸空间，m 是 T 上取值于 F 的向量测度，标量意义以 \(\mu\) 为基；对每个 \(z^{\prime} \in F^{\prime}\)，记 \(z^{\prime} \circ m = g_{z^{\prime}} \cdot \mu\)。为使 b) 的条件成立，必要且充分的是：对 T 的每个紧子集 K、每个等度连续子集 \(H^{\prime}\)（\(F^{\prime}\) 的）以及每个 \(\varepsilon > 0\)，存在 \(\delta > 0\)，使得关系式 \(A \subset K\) 与 \(\mu^{*}(A) \leqslant \delta\) 蕴含 \(\int_{A}^{*}|g_{z^{\prime}}| d\mu \leqslant \varepsilon\)，对每个 \(z^{\prime} \in H^{\prime}\) 成立（用上面的习题 11 与第 V 章 §5 习题 15）；这时称 m 关于 \(\mu\) 绝对连续（对 F 的原拓扑）。每个取值于半自反空间且标量意义以 \(\mu\) 为基的向量测度都关于 \(\mu\) 绝对连续（命题 3 的推论 2）。每个取值于 Banach 空间 \(\mathbf{F}\) 的可优超向量测度 m 都关于 \(|\mathbf{m}|\) 绝对连续（对 \(\mathbf{F}\) 的原拓扑）。
\end{enumerate}

¶ 13) 设 G 是 Hausdorff 桶式空间，F = G' 是其对偶；若当 F 赋以一个与 F 和 G 之间的对偶相容的拓扑时，m 是 T 上取值于 F 的向量测度，则当 F 赋以 G 的对偶的强拓扑时，m 仍是向量测度。

假设 m 标量意义以 \(\mu\) 为基；则 m 关于 \(\mu\) 绝对连续，当 F 赋以拓扑 \(\tau(\mathrm{F},\mathrm{G})\) 时（习题 12 c)）。为使当 F 赋以强拓扑时 m 关于 \(\mu\) 绝对连续，必要且充分的是：对 T 的每个紧子集 K 以及每个序列 \((\mathrm{A}_{n})\)（由 K 的 \(\mu\)-可测子集组成），在 \(\mathcal{L}^{1}(\mu)\) 中，由绝对值 \(\leqslant1\)、且在由序列 \(A_{n}\) 生成的集环（第 IV 章 §4，第 9 号）上的阶梯函数所成之集的闭包在 m 下的像 B，都包含一个对 \(G'\) 的强拓扑稠密的可数子集。（为看出条件是充分的，化归到 T = K 是 Stone 空间的情形（第 IV 章 §4，习题 10）。然后用反证法论证：若 \((\mathrm{A}_{n})\) 是由 K 的既开又闭的子集组成的序列，则在过渡到 K 的一个可度量化商空间 \(K_{1}\)（习题 10 a)）之后，可以假设 \(K_{1}\) 上与 \(\varphi_{A_{n}}\) 相对应的连续函数在 \(\mathcal{C}(\mathrm{K}_{1})\) 中组成一个完全集。另一方面，考虑 G 对与 B 正交的子空间的商，可以假设 B 在 \(F = G'\) 中对拓扑 \(\sigma(\mathrm{G}',\mathrm{G})\) 是完全的；若 H 是由 B 生成的 \(F = G'\) 的子空间，则假设蕴含 G 的每个有界子集 C 都对 \(\sigma(\mathrm{G},\mathrm{H})\) 预紧且可度量化。因此 G 的点的每个有界序列 \((\mathbf{z}_{n})\) 都有一个对 \(\sigma(\mathrm{G},\mathrm{H})\) 的 Cauchy 子序列；试证：若 \((\mathbf{z}_{n})\) 是对 \(\sigma(\mathrm{G},\mathrm{H})\) 的 Cauchy 序列，且记 \(z_{n} \circ m = g_{z_{n}} \cdot \mu\)，则 \(g_{z_{n}}\) 的类所成的序列在 \(L^{1}(\mu)\) 中对拓扑 \(\sigma(\mathrm{L}^{1},\mathrm{L}^{\infty})\) 是 Cauchy 序列；最后，利用第 V 章 §5 的习题 15 与 16 由此得出矛盾。）

\begin{enumerate}
\def\labelenumi{\arabic{enumi})}
\setcounter{enumi}{13}
\tightlist
\item
  \begin{enumerate}
  \def\labelenumii{\alph{enumii})}
  \tightlist
  \item
    设 G 是 Banach 空间，\(F = G'\) 是其强对偶，g 是 T 到 F 中的映射，属于空间 \(\Lambda_{G'}^{p}\)（§1，习题 16）。试证：若 p \textgreater{} 1，则测度 \(g \cdot \mu\) 关于 \(\mu\) 绝对连续（对 F 的强拓扑）。
  \end{enumerate}
\end{enumerate}

\begin{enumerate}
\def\labelenumi{\alph{enumi})}
\setcounter{enumi}{1}
\tightlist
\item
  取 G 为 Banach 空间 \(\mathrm{L}^{1}(\mathbf{N})\)，于是 \(\mathrm{G}^{\prime}=\mathrm{L}^{\infty}(\mathbf{N})\)；若 f 是 §1 习题 7a) 中定义的函数（视为取值于 \(G^{\prime}\) 中），则 f 是标量良可积的，但测度 \(f\cdot\mu\) 不关于 \(\mu\) 绝对连续（对 \(G^{\prime}\) 的强拓扑）。
\end{enumerate}

¶ 15) a) 设 f 是 T 到拟完备 Hausdorff 局部凸空间 F 中的标量局部可积映射。则 \(m = f \cdot \mu\) 是取值于 \(F'^{*}\) 的向量测度（对拓扑 \(\sigma(F'^{*}, F')\)）。试证：若这一测度是绝对连续的（对在 \(F'\) 的等度连续子集上一致收敛的拓扑），且此外 f 是 \(\mu\)-可测的（对 F 的原拓扑），则 m 是取值于 F 的向量测度。（利用 §1 第 2 号命题 8。）若 F 是 Banach 空间，且对某个 p \textgreater{} 1，\(\langle z', f \rangle\) 属于 \(\overline{\mathcal{L}}^{p}(\mu)\)（对一切 \(z' \in F'\)），则测度 m 是绝对连续的（参见习题 14 a））。

\begin{enumerate}
\def\labelenumi{\alph{enumi})}
\setcounter{enumi}{1}
\tightlist
\item
  设 \(\mathrm{I} = [0,1]\)，\(f\) 是定义在 \(\mathrm{I} \times \mathrm{I}\) 上（利用连续统假设）的函数，使得对每个 \(t \in \mathrm{I}\)，\(s \mapsto f(s,t)\) 是一个可数集的特征函数，而对每个 \(s \in \mathrm{I}\)，\(t \mapsto f(s,t)\) 是一个可数集的余集的特征函数（第 V 章 §8，习题 7 c)）。用 \(\mathrm{I}_0\) 表示赋有离散拓扑的区间 I，用 F 表示 I 上有界函数的 Banach 空间 \(\mathcal{B}(\mathrm{I}) = \mathrm{L}^\infty(\mathrm{I})\)；可把 F 与空间 \(\mathcal{C}(\widetilde{\mathrm{I}}_0)\)——即 \(\mathrm{I}_0\) 的 Stone-Čech 紧化上的连续函数所成的空间——等同起来。对每个 \(s \in \mathrm{I}\)，用 \(\mathbf{g}(s)\) 表示 F 的元素 \(t \mapsto f(s,t)\)；试证 \(\mathbf{g}\) 对 I 上的 Lebesgue 测度 \(\mu\) 是标量可积的；但 \(\int \mathbf{g} d\mu\) 不属于 F。更确切地说，\(\mathbf{g} \cdot \mu\) 是取值于 \(\mathrm{F}''\) 的向量测度，等于 \(\mathbf{c}\mu\)，其中 \(\mathbf{c}\) 是与 \(\widetilde{\mathrm{I}}_0 - \mathrm{I}_0 = \mathrm{E}\) 的特征函数等同的常向量（参见习题 12 a））。（把 \(\widetilde{\mathrm{I}}_0\) 上的每个正测度分解为一个原子测度与一个弥散测度之和（第 V 章 §5，第 10 号，命题 15），并注意弥散测度由 E 承载。）由此推出：不存在任何 \(\mu\)-可测函数 \(\mathbf{h}\)（取值于 \(\mathbf{F}\) 的），使 \(\mathbf{h} - \mathbf{g}\) 是标量 \(\mu\)-可忽略的（利用 \(a\)））。
\end{enumerate}

\begin{enumerate}
\def\labelenumi{\arabic{enumi})}
\setcounter{enumi}{15}
\item
  考察映射 \(u_{m}\)——T = {[}0,1{]} 到一个可分 Hilbert 空间中的、§1 习题 1 中定义的映射——并令 \(f_{k} = u_{1} + \cdots + u_{k}\)；试证向量测度 \(f_{k} \cdot \mu\) 在 \(\mathcal{C}(T)\) 的每个有界子集上一致收敛于一个取值于 F 的向量测度 m。此外，m 可连续延拓到空间 \(\mathcal{L}^{2}(\mu)\)，且对每个函数 \(f \in L^{2}\)，有 \(\int f dm \in F\) 以及在 F 中 \(|\int f dm| \leqslant N_{2}(f)\)；特别地，m 标量意义以 \(\mu\) 为基，且对强拓扑关于 \(\mu\) 绝对连续（习题 12 c)）。然而，对 T 上任何正测度 \(\nu\)，m 都不以 \(\nu\) 为基，从而（第 5 号，定理 1 的推论 4）更不是可优超的。（利用第 V 章 §5，第 7 号定理 3，化归到 \(\nu\) 以 \(\mu\) 为基的情形。）
\item
    设 \(\mu\) 是 \(T = [0,1]\) 上的 Lebesgue 测度；映射 \(f \mapsto f \cdot \mu\)（\(\mathcal{C}(T)\) 到 \(\mathcal{M}(T)\) 中的）对 \(\mathcal{M}(T)\) 上的强拓扑连续，因而是取值于该 Banach 空间的向量测度 \(\mathbf{m}\)。试证 \(\mathbf{m}\) 标量意义以 \(\mu\) 为基，且对强拓扑关于 \(\mu\) 绝对连续，但 \(\mathbf{m}\) 不以 \(\mu\) 为基（对 \(\mathcal{M}(T)\) 上的强拓扑）。（注意 \(\mathbf{m}\) 以 \(\mu\) 为基（对淡拓扑 \(\sigma(\mathcal{M}(T), \mathcal{C}(T))\) 而言），并由此推出：若有 \(\mathbf{m} = \mathbf{g} \cdot \mu\)，其中 \(\mathbf{g}\) 标量 \(\mu\)-可积（对 \(\mathcal{M}(T)\) 上的强拓扑），则几乎处处必有 \(\mathbf{g}(t) = \varepsilon_t\)。试证这导致矛盾，办法是注意：对每个有界数值函数 \(\theta\)（\(T\) 上的），线性形式 \(\mathbf{z}' : \lambda \mapsto \sum_{t \in T} \theta(t) \lambda(\{t\})\)
\end{enumerate}

在 \(\mathcal{M}(\mathrm{T})\) 上连续，但 \(\langle \mathbf{g},\mathbf{z}'\rangle\) 不一定 \(\mu\)-可测。）

¶ 18) a) 设 T 是局部紧空间，\((K_{\alpha})\) 是 T 的由紧集组成的局部有限覆盖。设 \(\mu\) 是 T 上的正测度，\(\mu_{\alpha}\) 是 \(\mu\) 在 \(K_{\alpha}\) 上诱导的测度。假设每个空间 \(L^{\infty}(K_{\alpha}, \mu_{\alpha})\) 都具有提升性质。试证 \(L^{\infty}(T, \mu)\) 具有提升性质。

\begin{enumerate}
\def\labelenumi{\alph{enumi})}
\setcounter{enumi}{1}
\tightlist
\item
  设 K 是可度量化紧空间，\(\mu\) 是 K 上的正测度，\((\varpi_n)\) 是由 K 的分成可积集的有限分割组成的基本序列（第 IV 章 §5，习题 17）。对 K 的每个可积子集 A 以及每个元素 \(\widetilde{f} \in \mathrm{L}^1(\mu)\)，令 \(\lambda_A(\widetilde{f}) = 0\)（若 \(\mu(A) = 0\)），否则令 \(\lambda_A(\widetilde{f}) = (\mu(A))^{-1} \int_{A} f d\mu\)；对每个 \(n\)，令 \(\rho_n(\widetilde{f}) = \sum_k \lambda_{A_k}(\widetilde{f}) \varphi_{A_k}\)，如果
\end{enumerate}

\(\varpi_{n} = (\mathrm{A}_{k})\)。试证序列 \((\rho_{n}(\tilde{f}))\)（\(\mu\)-可积函数的序列）几乎处处收敛于 \(f\)。（先考虑 \(f\) 是连续函数的情形。给定任意数 \(a > 0\)，再证明：所有属于至少一个分割 \(\varpi_{n}\) 且满足 \(a \cdot \mu(A) \leqslant \int_{A} f d\mu\) 的集 A 的并 B 是可测的，且 \(\mu(B) \leqslant a^{-1} \int f d\mu\)。最后对 \(f\)，用一列连续函数在 \(\mathcal{L}^{1}\) 中逼近它。）

\begin{enumerate}
\def\labelenumi{\alph{enumi})}
\setcounter{enumi}{2}
\item
  试证每个可度量化紧空间都具有提升性质。（设 \(\mathfrak{U}\) 是 \(\mathbf{N}\) 上比 Fréchet 滤子更细的超滤子；试证 \(\rho(\widetilde{f}) = \lim_{\mathfrak{U}} \rho_n(\widetilde{f})\) 是 \(\mathrm{L}^{\infty}\) 的一个提升。）
\item
  由 \(a\) 与 \(c\) 推出：每个可度量化局部紧空间都具有提升性质（对任何正测度均是）。
\end{enumerate}

\begin{enumerate}
\def\labelenumi{\arabic{enumi})}
\setcounter{enumi}{18}
\tightlist
\item
  设 \(\mu\) 是 T 上的测度，使得 Banach 空间 \(L^1 (\mu)\) 可分。试证：\(L^1 (\mu)\) 到强对偶 \(F'\)（任意 Banach 空间 \(F\) 的强对偶）中的每个连续线性映射，都可以由形如 \(g\mapsto \int g\mathbf{f}d\mu\) 的映射通过过渡到商而得到，其中 \(\mathbf{f}\in \mathcal{L}_{\mathrm{F}_s'}^\infty\)。（利用 TVS，IV，§2 习题 19，化归到 Dunford-Pettis 定理。）
\end{enumerate}

¶ 20) 设 F 是可分 Fréchet 空间，F' 是其对偶，μ 是 T 上的正测度，m 是 T 上取值于弱对偶 \(F'_s\) 的测度，标量意义以 μ 为基。为使 m 以 \(\mu\) 为基，必要且充分的是下列条件成立：对每个 \(\varepsilon > 0\) 以及每个紧子集 \(K\)（\(T\) 的），存在紧子集 \(K_1 \subset K\)，使 \(\mu(K - K_1) \leqslant \varepsilon\)，且使得在 \(m\) 之下，由所有 \(\mu\)-可测函数 \(g\)（有界、支集含于 \(K_1\) 且满足 \(\int |g| d\mu \leqslant 1\) 的）组成的集的像，是 \(F'\) 的等度连续子集。（回忆：\(\int g dm \in F'\) 对每个有界且具有紧支集的 \(\mu\)-可测函数 \(g\) 都成立；参见命题 3 的推论 2。为证条件必要，用 §1 第 5 号命题 13 与 §1 第 2 号命题 5。为证条件充分，先考虑 \(T\) 紧的情形；应用第 5 号定理 1 的推论 3，先证明存在 \(T\) 分解为一个可忽略集 \(N\) 与一列紧集（\(K_n\)）的分割，以及一个可测映射 \(g\)（\(T\) 到 \(F'_s\) 中的），使得 \(\int f dm = \int g f d\mu\) 对每个函数 \(f \in \mathcal{L}^\infty(\mu)\) 都成立，只要该函数在有限个集 \(K_n\) 的并的余集上为零；为证明 \(m = g \cdot \mu\)，用第 V 章 §5 的习题 16 b) 与 20 a)。最后，为过渡到 \(T\) 是任意局部紧空间的情形，用第 IV 章 §5，第 9 号命题 14。）

¶ 21) 设 F 是 Banach 空间，u 是 Banach 空间 \(\mathrm{L}_{\mathrm{F}}^{p}(\mathrm{T},\mu)\) 上的连续线性形式；用 q 表示 p 的共轭指数，用 E 表示 \(\mu\)-可积集上的数值阶梯函数所成的空间。对 \(f \in L^{p}\) 与 \(z \in F\)，可写 \(u(fz) = \langle z, m(f) \rangle\)，其中 m 是 \(L^{p}\) 到 \(F'\) 中的连续线性映射，满足 \(|m(f)| \leqslant \|u\| \cdot N_{p}(f)\)。

\begin{enumerate}
\def\labelenumi{\alph{enumi})}
\tightlist
\item
  设 \((\mathrm{A}_i)_{1\leqslant i\leqslant n}\) 是由两两不交且非可忽略的 \(\mu\)-可积集组成的有限序列。试证
\end{enumerate}

\[
\sum_ {i} \left(| \mathbf {m} (\varphi_ {\mathrm{A} _ {i}}) | ^ {q} / (\mu (\mathrm{A} _ {i})) ^ {q - 1}\right) \leqslant \| u \| ^ {q}.
\]

（对每个向量系统 \((\mathbf{a}_i)_{1\leqslant i\leqslant n}\)（F 中的向量，满足 \(|\mathbf{a}_i| = 1\)），考虑线性形式 \((\xi_i)\mapsto u\big(\sum_i\xi_i\mathbf{a}_i\varphi_{\mathrm{A}_i}\big)\)（\(\mathbf{R}^n\) 上的），并对一个有限支集的测度应用第 V 章 §5，第 8 号定理 4。）由此推出：若 \(\mathrm{A} = \bigcup_{i}\mathrm{A}_{i}\)，则 \(\sum_{i}|\mathbf{m}(\varphi_{\mathrm{A}_i})|\leqslant \| u\| (\mu (\mathrm{A}))^{1 / p}\)（应用 Minkowski 不等式）。

\begin{enumerate}
\def\labelenumi{\alph{enumi})}
\setcounter{enumi}{1}
\tightlist
\item
  对每个正函数 \(f \in \mathcal{E}\)，令 \(\nu(f) = \sup\left(\sum_{i} |\mathbf{m}(f\varphi_{\mathrm{A}_i})|\right)\)，其中 \((\mathrm{A}_i)\)
\end{enumerate}

取遍由两两不交的 \(\mu\)-可积集组成的有限序列所成之集。试证 \(\nu\) 是某个正测度在 \(\mathcal{E}\) 上的限制——该正测度以 \(\mu\) 为基（仍记作 \(\nu\)），定义在 T 上——使得 \(|\mathbf{m}(f)| \leqslant \nu(f)\) 对每个正函数 \(f \in \mathcal{E}\) 都成立（利用 \(a\)））。

\begin{enumerate}
\def\labelenumi{\alph{enumi})}
\setcounter{enumi}{2}
\tightlist
\item
  试证：若 F 可分，则存在 T 到 F' 中的映射 g，标量 μ-可测，使得 \(u(\widetilde{\mathbf{f}})=\int\langle\mathbf{f},\mathbf{g}\rangle d\mu\) 对每个函数 \(f\in L_{F}^{p}\) 都成立，且 \(\|u\|=N_{q}(\mathbf{g})\)。（为证最后的等式，用 §1 第 5 号命题 13 与 §1 习题 13。）d) 由 c) 推出：若 F 是自反可分 Banach 空间，则空间 \(L_{F}^{p}\) 对 \(1<p<+\infty\) 自反（应用 §1 第 5 号命题 12）。
\end{enumerate}

¶ 22) 设 F 是 Hausdorff 局部凸空间，S 是 F 的子集所成之集，这些子集是凸的、均衡的且对 \(\sigma(F, F')\) 紧，而 \(S'\) 是 \(F'\) 的子集所成之集，这些子集是凸的、等度连续的、均衡的且对 \(\sigma(F', F'')\) 紧。

\begin{enumerate}
\def\labelenumi{\alph{enumi})}
\item
  假设 \(\mathfrak{S}\) 中每个集都对 \(\mathfrak{S}'\)-拓扑预紧。试证：每个连续线性映射 \(u\)——\(F\) 到拟完备 Hausdorff 局部凸空间 \(G\) 中的连续线性映射——只要把 \(F\) 的有界子集变为对 \(\sigma(G, G')\) 相对紧的子集，就把每个属于 \(\mathfrak{S}\) 的集变为 \(G\) 中对原拓扑相对紧的子集。（用 TVS，IV，§1 习题 4 与上面的习题 11。）逆命题（考虑 \(F\) 到其对 \(\mathfrak{S}'\)-拓扑的完备化空间中的典范映射）。
\item
  假设 F 是可度量化空间，或是可度量化空间的严格直接极限。试证：若每个对 \(\sigma(\mathrm{F},\mathrm{F}^{\prime})\) 的 Cauchy 序列都是对 \(S^{\prime}\)-拓扑的 Cauchy 序列，则 a) 的条件满足（利用 Šmulian 定理（TVS，IV，§5，习题 2 c）））。
\item
  假设 F 是拟桶式的（TVS，III，§4，习题 7），并用 \(\mathfrak{S}''\) 表示 \(\mathrm{F}''\) 的子集所成之集，这些子集是凸的、等度连续的、均衡的且对 \(\sigma(\mathrm{F}''\) , \(\mathrm{F}'''\) ) 紧。试证：若 \(\mathfrak{S}'\) 中每个集都对 \(\mathfrak{S}''\)-拓扑预紧，则 \(\mathfrak{S}\) 中每个集都对 \(\mathfrak{S}'\)-拓扑预紧。（用 TVS，IV，§1 习题 4，并注意 F 的原拓扑由 \(\mathrm{F}''\) 上的强拓扑诱导。）
\end{enumerate}

¶ 23) 设 T 是局部紧空间，\(\mu\) 是 T 上的正测度，F 是三个 Banach 空间 \(\overline{\mathcal{K}(T)}\)、\(L^1(\mu)\)、\(L^\infty(\mu)\) 之一。设 u 是 F 到拟完备空间 G 中的连续线性映射，它把 F 的单位球变为 G 中对 \(\sigma(G, G')\) 相对紧的子集。试证 u 把 F 的每个对 \(\sigma(F, F')\) 相对紧的子集变为 G 中对原拓扑相对紧的子集。（应用习题 22 c)，把情形 \(F = L^1(\mu)\) 化归到情形 \(F = L^\infty(\mu)\)；用第 II 章 §1 习题 13 把情形 \(L^\infty(\mu)\) 化归到情形 \(F = \mathcal{C}(S)\)，其中 S 是一个适当的紧空间。若 \(F = \overline{\mathcal{K}(T)}\)，则应用习题 22 b)，然后用第 V 章 §5 的习题 22 b) 与 15；注意对 \(\sigma(F, F')\) 的 Cauchy 序列在 T 上逐点收敛，从而对 T 上每个有界测度依测度收敛。）

¶ 24) 设 \(\mu\) 是 \(T\) 上的正测度，\(u\) 是 \(L^1(\mu)\) 到 Banach 空间 \(F\) 中的线性映射，把 \(L^1(\mu)\) 的单位球变为 \(F\) 中对 \(\sigma(F, F')\) 相对紧的子集。试证：存在 \(\mu\)-可测映射 \(f\)（\(T\) 到 \(F\) 中的），使 \(|\mathbf{f}(t)| \leqslant \|u\|\) 对一切 \(t \in T\) 成立，且使得

\[
u (\widetilde {g}) = \int \mathbf {f} g d \mu
\]

对每个函数 \(g \in \overline{\mathcal{L}}^1(\mu)\)（Dunford-Pettis-Phillips 定理）。（先借助第 IV 章 §5，第 9 号命题 14 化归到 T 紧的情形。此时 \(L^\infty(\mu) \subset L^1(\mu)\) 的单位球 B 对 \(\sigma(L^1, L^\infty)\) 相对紧（第 V 章 §5，习题 15）；利用习题 23，试证由 \(u(L^1)\) 生成的 F 的闭线性子空间是可分的。于是可化归到 F 可分且 \(\overline{u(L^1)} = F\) 的情形。此时在 \(u\) 之下，\(L^1\) 的单位球的像在 F 中的闭包 A 对 \(\sigma(F, F')\) 紧（第 IV 章 §7，习题 10）；试证它对该拓扑可度量化（TVS，III，§3，第 4 号，命题 6 的推论 2）。注意此时 A 可等同于一个赋范空间 G 的对偶的单位球，其中 G 是把 \(F'\) 赋以等于 \(A^\circ\) 的规范函数的范数而得到的。试证 G 可分，从而化归到应用 Dunford-Pettis 定理（定理 1 的推论 2）。

¶ 25) 设 \(\mu\) 是 T 上的正测度。

\begin{enumerate}
\def\labelenumi{\alph{enumi})}
\item
  设 \(\mathbf{f}\) 是标量 \(\mu\)-可测映射（\(\mathrm{T}\) 到 Banach 空间 \(\mathrm{F}\) 中的），使得 \(\mathbf{f}(\mathrm{T})\) 对 \(\sigma(\mathrm{F},\mathrm{F}')\) 相对紧。试证：存在 \(\mu\)-可测映射 \(\mathbf{g}\)（\(\mathrm{T}\) 到 \(\mathrm{F}\) 中的），使 \(\mathbf{f} - \mathbf{g}\) 是标量局部可忽略的。（考虑映射 \(\widetilde{h} \mapsto \int \mathbf{f}h d\mu\)（\(\mathrm{L}^1(\mu)\) 到 \(\mathrm{F}\) 中的），对它应用习题 24，并用第 IV 章 §7 习题 10。）
\item
  设 \(\mathbf{f}\) 是 T 到自反 Fréchet 空间 F 中的标量 \(\mu\)-可测映射。试证：存在 T 到 F 中的 \(\mu\)-可测映射 \(\mathbf{g}\)，使 \(\mathbf{f} - \mathbf{g}\) 是标量局部可忽略的（参见习题 15 b））。（用第 IV 章 §5，第 9 号命题 14，化归到 T 紧的情形。然后应用 §1 习题 23 c)，化归到 \(\mathbf{f}(T)\) 在 F 中有界的情形。最后把 F 嵌入可数个 Banach 空间的乘积，并用 a)。）
\item
  设 \(\mathbf{f}\) 是 \(\mathrm{T}\) 到 Fréchet 空间 \(\mathrm{F}\) 中的映射，它是 \(\mu\)-可测的（对拓扑 \(\sigma(\mathrm{F}, \mathrm{F}')\)）。试证 \(\mathbf{f}\) 是 \(\mu\)-可测的（对 \(\mathrm{F}\) 的原拓扑）。（化归到 \(\mathrm{T}\) 紧且 \(\mathbf{f}\) 对拓扑 \(\sigma(\mathrm{F}, \mathrm{F}')\) 连续的情形；然后像 \(b\) 那样结束证明。）
\end{enumerate}

¶ 26) 设 S、T 是两个紧空间，f 是定义在 S × T 上的有限数值函数。

\begin{enumerate}
\def\labelenumi{\alph{enumi})}
\item
  为使映射 \(s \mapsto f(s, \cdot)\) （S 到 \(\mathbf{R}^{\mathrm{T}}\) 中的）是 S 到空间 \(\mathcal{C}(\mathrm{T})\) （赋以拓扑 \(\sigma(\mathcal{C}(\mathrm{T}), \mathcal{M}(\mathrm{T}))\) ）中的连续映射，必须且只需 \(f\) 有界，且每个部分映射 \(f(s, \cdot), f(\cdot, t)\) ( \(s \in \mathrm{S}, t \in \mathrm{T}\) ) 都连续。（为证条件充分，先证明 S 在 \(s \mapsto f(s, \cdot)\) 下的像 M 对 \(\sigma(\mathcal{C}(\mathrm{T}), \mathcal{M}(\mathrm{T}))\) 是紧的。为此，注意 \(s \mapsto f(s, \cdot)\) 当 \(\mathcal{C}(\mathrm{T})\) 赋以 T 上逐点收敛的拓扑时是连续的；利用 Eberlein 定理（TVS，IV，§5，第 3 号，定理 1）化为证明每个序列 \((f(s_n, \cdot))\) 在 \(\mathcal{C}(\mathrm{T})\) 中对 \(\sigma(\mathcal{C}(\mathrm{T}), \mathcal{M}(\mathrm{T}))\) 有一个聚点。通过考虑 T 的一个适当商空间（参见习题 10 a)）化为 T 可度量化的情形，并注意，在 \(\mathcal{C}(\mathrm{T})\) 中对 T 上逐点收敛拓扑相对紧的子集上，该拓扑与在 T 的稠密子集上逐点收敛的拓扑相同。由此得到 \((f(s_n, \cdot))\) 的一个在 T 上逐点收敛的子序列，并应用 Lebesgue 定理。最后注意，在 M 上拓扑 \(\sigma(\mathcal{C}(\mathrm{T}), \mathcal{M}(\mathrm{T}))\) 与逐点收敛的拓扑相同。）
\item
  假设每个部分映射 \(f(s, \cdot), f(\cdot, t)\) 都连续 ( \(s \in S\) , \(t \in T\) )。证明：对每个正测度 \(\mu\) ( \(S\) 上的)及每个 \(\varepsilon > 0\) ，存在紧子集 \(K \subset S\) ，使得 \(\mu(S - K) \leqslant \varepsilon\) 且 \(f\) 在 \(K \times T\) 上的限制连续。（化为 \(f\) 有界的情形；应用 a) 与习题 25 c)。）
\item
  在与 b) 相同的假设下，证明 f 对 S × T 上每个测度 \(\nu\) 都可测。（对 \(\nu\) 在 S × T 到 S 上的射影下的像应用 b)。）
\end{enumerate}

\begin{enumerate}
\def\labelenumi{\arabic{enumi})}
\setcounter{enumi}{26}
\tightlist
\item
  设 \(\mathbf{m}\) 是 \(\mathbf{T}\) 上的向量测度，取值于一个 Hausdorff 局部凸空间 \(\mathbf{F}\) 。
\end{enumerate}

\begin{enumerate}
\def\labelenumi{\alph{enumi})}
\item
  证明：若 \(\mathbf{m}\) 的支集有限，则 \(\mathbf{m} = \sum_{i=1}^{n} \mathbf{c}_i \varepsilon_{a_i}\) ，其中 \(\mathbf{c}_i \in \mathbf{F}\) 。
\item
  若 \(\mathbf{F}\) 是 Banach 空间，且 \(\mathbf{m}\) 对紧收敛拓扑连续，证明 \(\mathbf{m}\) 具有紧支集。
\item
  给出一个测度的例子：取值于 \(\mathbf{R}^{\mathbf{N}}\) ，支集非紧，且对紧收敛拓扑连续。
\end{enumerate}

\setcounter{exercisegroup}{3}
\edef\CurrentExerciseGroup{\arabic{exercisegroup}}
\section*{\texorpdfstring{\S\CurrentExerciseGroup}{第{\CurrentExerciseGroup}节}}
\phantomsection
\addcontentsline{toc}{section}{第{\CurrentExerciseGroup}节习题}

\begin{enumerate}
\def\labelenumi{\arabic{enumi})}
\item
  设第 1 号的定理 1 之假设成立。设 \(h\) 是局部 \(\mu\) -可积函数，使得 \(p\) 是 \((h \cdot \mu)\) -逆紧的。证明函数 \(b \mapsto g(b) = \int h(t) d\lambda_b(t)\) ——它对测度 \(\nu = p(\mu)\) 而言在 B 中局部几乎处处有定义——满足 \(p(h \cdot \mu) = g \cdot \nu\) 。
\item
  设 B 是局部紧空间，\((\nu_{\iota})_{\iota\in I}\) 是 B 上全体正测度构成的族。设 T 是乘积空间 \(I\times B\) （其中 I 赋有离散拓扑），设 \(\nu\) 是 T 上这样的测度：它在 \(\{\iota\}\times B\) 上的限制是 B 上测度 \(\nu_{\iota}\) 的典范像（对一切 \(\iota\in I\) ）。设 p 是 T 到 B 上的射影；证明：若 B 中存在不可数的紧子集，则 \(\nu\) 在 p 下不容许伪像测度。（证明对这样的测度，B 的每一点都将有测度 \textgreater0 。）
\item
  给出如下例子：正测度 \(\mu\) （在一个波兰局部紧空间 \(T\) 上的）与连续映射 \(p\) ( \(T\) 到一个波兰局部紧空间 \(B\) 中的)，使得 \(p\) 不是 \(\mu\) -逆紧的。
\item
  设 \(\mathrm{T}\) 是区间 {[}0,1{]} ( \(\mathbf{R}\) 中的)，\(\mu\) 是 \(\mathrm{T}\) 上的 Lebesgue 测度。设 \(\mathrm{R}\{x,y\}\) 为等价关系 \(x - y \in \mathbf{Q}\) （在 \(\mathrm{T}\) 中）。证明 \(\mathrm{R}\) 不是 \(\mu\) -可测的（应用第 4 号定理 3），但 \(\mathrm{R}\) 在 \(\mathrm{T} \times \mathrm{T}\) 中的图像对乘积测度 \(\mu \otimes \mu\) 是可忽略的。
\item
  设 T 是 Cantor 三分集 \(K \subset [0,1]\) 与区间 I = {[}1,2{]}（R 中的）的并，设 \(\mu\) 是 Lebesgue 测度在紧空间 T 上诱导的测度。设 P 是 I 的一个具有连续统势的不可测子集（第 IV 章 §4，习题 8），设 \(\psi\) 是 K 到 P 上的双射。考虑 T 中的等价关系 R：每个不属于 \(K \cup P\) 的点 x 自成一个等价类，而点 \(y \in K\) 的类由 y 与 \(\psi(y)\) 组成。证明 R 是 \(\mu\) -可测的，但 K 对 R 的饱和化不是 \(\mu\) -可测的。
\item
  设 T 是波兰局部紧空间，\(\mu\) 是 T 上的正测度，R 是 T 中的 \(\mu\) -可测等价关系，p 是 T 到波兰局部紧空间 B 中的 \(\mu\) -可测映射，使得 \(p(x) = p(y)\) 等价于 \(R\{x, y\}\) ，\(\nu\) 是 \(\mu\) 在 p 下的伪像测度，而 \(b \mapsto \lambda_b\) 是 \(\mu\) 按关系 R 的一种分解。对每个 \(b \in B\) (满足 \(\lambda_b \neq 0\) 的)，设 \(C(b)\) 是承载 \(\lambda_b\) 的模 R 等价类；证明 \(\varphi_{C(b)} \cdot \mu\) 与 \(\lambda_b\) 成比例。给出一个例子，其中 \(\varphi_{C(b)} \cdot \mu = 0\) 对每个 \(b \in B\) 成立。
\item
  设 T 是波兰局部紧空间，\(\mu\) 是 T 上的有界正测度，R 是 T 中的 \(\mu\) -可测等价关系。设 \(\Omega\) 是 T 上测度空间 \(\mathcal{M}(\mathrm{T})\) 中由满足 \(\lambda \geqslant 0\) 、总质量 \(\leqslant 1\) 且非零的测度组成的子集；\(\Omega\) 赋有淡拓扑（第 III 章 §1，第 9 号，命题 15 的推论 2）时是局部紧的。证明存在唯一的正测度 \(\rho\) （ \(\Omega\) 上的），使得： \(1^{\circ} \mu = \int_{\Omega} \lambda d\rho(\lambda)\) ； \(2^{\circ} \rho\) 集中于一个子集 \(B_{0}\) （ \(\Omega\) 中由总质量为 1、由模 R 的类承载的测度组成的子集）上，且 \(B_{0}\) 的两个不同元素由不同的类承载。（考虑分解 \(b \mapsto \lambda_{b}\) （ \(\mu\) 按关系 R 的一种分解），使所有 \(\lambda_{b}\) 都具有总质量 1，并考虑 B 在映射 \(b \mapsto \lambda_{b}\) （ B 到 \(\Omega\) 中的）下的像；利用定理 4。）
\item
  设 T 是 R 中的区间 \([0,2]\) ，\(\mu\) 是 T 上的 Lebesgue 测度。设 A 是含于 Cantor 集 \(K \subset [0,1]\) 的非 Borel 集（参见 GT，IV，§2，第 4 号，例以及 IX，§6，习题 6）。设 S 是 T 中的等价关系，其等价类是集合 \(\{x\}\) (对 \(x \notin A \cup (A + 1)\) )与集合 \(\{x, x + 1\}\) (对 \(x \in A\) )。证明 S 是 \(\mu\) -可测的，但不存在 S 的 Borel 截面。
\item
  设 \(\mathbf{K}\) 是可度量化紧空间，\(f\) 是 \(\mathbf{K}\) 到 Hausdorff 拓扑空间 \(\mathbf{E}\) 中的连续映射，\(k\) 是任意整数 \(> 1\) 。用 \(\mathrm{B}_k\) 表示 \(\mathbf{E}\) 中由这样的 \(y\) 组成的子集：\(\stackrel{-1}{f(y)}\) 至少含 \(k\) 个不同的点；证明 \(\mathrm{A}_k = \stackrel{-1}{f(\mathrm{B}_k)}\) 是 \(\mathbf{K}\) 中的 Borel 集（对每个整数 \(n > 0\) ，用 \(\mathrm{B}_{kn}\) 表示 \(\mathbf{E}\) 中由这样的 \(y\) 组成的子集：\(\stackrel{-1}{f(y)}\) 至少含 \(k\) 个彼此距离均为 \(\geqslant 1/n\) 的点；证明 \(\mathrm{B}_{kn}\) 是闭集）。
\end{enumerate}

¶ 10) 设 K 是可度量化紧空间，μ 是 K 上的正测度，f 是 K 到 Hausdorff 拓扑空间 E 中的 μ-可测映射。用 I（相应地 U）表示 E 中由这样的 y 组成的子集：\(\stackrel{-1}{f(y)}\) 是无限集（相应地不可数集）。

\begin{enumerate}
\def\labelenumi{\alph{enumi})}
\item
  证明 K 中存在 \(\mu\) -可测集 H，使得 \(f(\mathrm{H}) = \mathrm{I}\) ，且对每个 \(y \in \mathrm{E}\) ，\(\stackrel{-1}{f(y)} \cap \mathbb{C}\mathrm{H}\) 是有限集。（设 \((\mathrm{K}_n)\) 是 K 中紧子集的递增序列，使得 \(f\) 在每个 \(\mathrm{K}_n\) 上的限制连续，且诸 \(\mathrm{K}_n\) 之并 F 的余集 N 是 \(\mu\) -可忽略的；对每个 \(k > 1\) ，设 \(\mathrm{B}_k\) 是 E 中由这样的 \(y\) 组成的子集：\(\mathrm{F} \cap \stackrel{-1}{f(y)}\) 至少含 \(k\) 个不同的点，并设 \(\mathrm{A}_k = \mathrm{F} \cap \stackrel{-1}{f(\mathrm{B}_k)}\) ；取 H 为集合 \(\mathrm{A} = \bigcap_k \mathrm{A}_k\) 与集合 \(\mathrm{N} \cap \stackrel{-1}{f(\mathrm{I})}\) 之并，并利用习题 9。）
\item
  设 \(\mathfrak{F}\) 是由这样的 \(\mu\) -可测集 \(M \subset H\) 组成的集：\(M \cap f(y)\) 对每个 \(y \in E\) 都是有限集；设 \(\alpha\) 是测度 \(\mu(M)\) ( \(M \in \mathfrak{F}\) )的上确界；存在集合的递增序列 \((\mathrm{M}_n)\) （ \(\mathfrak{F}\) 中的），使得 \(\lim_{n\to \infty}\mu (\mathrm{M}_n) = \alpha\) 。设 \(\mathrm{P} = \bigcup_{n}\mathrm{M}_{n}\)
\end{enumerate}

并设 \(S\) 是 \(H \cap \mathbb{C}P\) 关于等价关系 \(f(x) = f(y)\) 的可测截面（定理 3）；证明 \(\mu(S) = 0\) 。

\begin{enumerate}
\def\labelenumi{\alph{enumi})}
\setcounter{enumi}{2}
\tightlist
\item
  证明存在 \(\mu\) -可忽略集 \(Z \subset K\) ，使得 \(f(Z) = U\) ，且对每个 \(\varepsilon > 0\) ，存在 \(\mu\) -可测集 \(L \subset K\) ，使得 \(\mu(L) \leqslant \varepsilon\) 且 \(f(L) = I\) （注意 \(f(H \cap \mathbb{C}M_n) = I\) 对每个 \(n\) 成立，且 \(U \subset f(S)\) ）。
\end{enumerate}

¶ 11) 设 K 是可度量化紧空间，μ 是 K 上的正测度，f 是 K 到局部紧空间 T 中的 μ-可测映射，ν 是 T 上的正测度；I 与 U 的含义同习题 10。

\begin{enumerate}
\def\labelenumi{\alph{enumi})}
\item
  假设对每个 \(\mu\) -可忽略集 \(N \subset K\) ，\(f(N)\) 都是 \(\nu\) -可忽略的。则 \(U\) 是 \(\nu\) -可忽略的，且对每个 \(\mu\) -可测子集 \(A \subset K\) ，\(f(A)\) 是 \(\nu\) -可测的。
\item
  为使 I 是 \(\nu\) -可忽略的且在 \(f\) 下每个 \(\mu\) -可忽略集的像是 \(\nu\) -可忽略的，必须且只需 \(f\) 满足以下条件：对每个 \(\varepsilon > 0\) ，存在 \(\delta > 0\) ，使得对每个子集 \(A \subset K\) ( \(\mu\) -可测且满足 \(\mu(A) \leqslant \delta\) 的)，\(f(A)\) 是 \(\nu\) -可测的且 \(\nu(f(A)) \leqslant \varepsilon\) 。（为证条件必要，用反证法，考虑递减序列 \((A_n)\) （由 \(\mu\) -可测集组成），其交是 \(\mu\) -可忽略的，但 \(\nu(f(A_n)) \geqslant \alpha > 0\) ；然后取诸 \(f(A_n)\) 与 \(\mathbb{C}\mathrm{I}\) 的交。）
\end{enumerate}

¶ 12) 设 E 是在 R 的区间 \([-1,+1]\) 上赋以 GT，IX，§2 的习题 13 a) 中定义的拓扑 T 而得到的非可度量化紧空间；设 \(\varphi\) 是映射 \(x \mapsto |x|\) （ E 到 I = {[}0,1{]} 上的，I 赋以 R 的拓扑所诱导的拓扑）。

\begin{enumerate}
\def\labelenumi{\alph{enumi})}
\item
  证明 \(\varphi\) 连续，并且，若用 \(\mathfrak{F}\) 表示 \(E\) 的形如 \(\overline{\varphi}(A)\) 的子集组成的集（其中 \(A\) 是 \(I\) 的 Borel 子集），则 \(E\) 的 Borel 子集是这样的子集 \(B\) ：存在 \(M \in \mathfrak{F}\) ，使得 \(B \cap \mathbb{C}M\) 与 \(M \cap \mathbb{C}B\) 均可数（注意 \(E\) 的每个开集都属于 \(\mathfrak{F}\) ）。
\item
  证明数值函数 \(f\) （定义在 \(\mathbf{E}\) 上的）为连续（对 \(\mathcal{T}\) ），必须且只需它是正则函数，且 \(f(x-) = f((-x)+)\) 对每个 \(x \in \mathbf{E}\) 成立。于是定义正测度 \(\mu\) （ \(\mathbf{E}\) 上的）如下：令 \(\int f d\mu = \int_0^1 f(x) dx\) 。证明 \(\mu\) 在 \(\varphi\) 下的像是 \(\mathbf{I}\) 上的 Lebesgue 测度，且 \(\mu\) -可忽略集正是 \([-1, +1]\) 上对 Lebesgue 测度可忽略的子集。
\item
  由 \(a)\) 与 \(b)\) 推出：不存在 \(\mu\) -可测截面，对应于 \(\mu\) -可测等价关系 \(\varphi(x) = \varphi(y)\) （在 \(E\) 中）。
\end{enumerate}

\chapter*{历史注记}
\phantomsection
\addcontentsline{toc}{chapter}{历史注记}

（注意 --- 罗马数字指本注记末尾的参考文献。）

随着 19 世纪“向量分析”的发展，积分向量值函数成为常见之事，但只要所涉及的只是取值于有限维空间的函数，这一运算并不构成问题。只是到了 Hilbert 的谱理论，才遇到自然导向更一般的积分概念的运算：这一理论实际上导致把 Hilbert 空间 H 上的每个连续 Hermite 型 \(\Phi(x,y)\) 对应于一个正交射影族 \((E(\lambda))_{\lambda\in\mathbf{R}}\) ，它具有这样的性质：对 H 的每一对向量 \((x,y)\) ，函数 \(\lambda\mapsto(E(\lambda)x|y)\) 是有界变差的，且 \(\Phi(x,y)=\int\lambda d(E(\lambda)x|y)\) ；若把 \(\Phi\) 与使 \(\Phi(x,y)=(Ax|y)\) 成立的 Hermite 算子 A 相联系，人们便很想把前一公式写成 \(A=\int\lambda dE(\lambda)\) 。但只是从大约 1935 年起，在 Bochner 引进了取值于 Banach 空间的函数的（“强”）积分之后，人们才开始致力于定义向量值函数的积分（或关于向量测度的积分），使之能够合法地写出诸如前一公式那样的公式。这一推广主要是由 Gelfand（III）、Dunford 与 Pettis（IV）和（V）完成的；他们的结果是对 Banach 空间叙述的，但可毫无困难地推广到更一般的局部凸空间。

把一个体积分解成“切片”、并把该体积上的积分化为每个切片上的积分、随后做一次单独积分的想法，在分析学中自无穷小演算的发端起即已被使用（Cavalieri 的“不可分元演算”不过是这一原理的最初轮廓，它甚至可以追溯到 Archimedes（参见第 IV 卷第 I--III 章的历史注记））。但在经典的应用中，“切片”总是具有非常特殊、非常正则的性质（最常见的是解析曲面的开子集，解析地依赖于一个参数）；在没有一般的积分理论的情况下，也很难有别的情形。测度分解的一般问题由 von Neumann 于 1932 年在遍历理论的关联中提出并解决（I）；大约在同一时间（并且独立地），Kolmogoroff 在为概率论奠定公理化基础时，被引向以一般的方式定义“条件概率”概念并证明其存在性，这个问题本质上等价于测度的分解问题（II）。

\chapter*{参考文献}
\phantomsection
\addcontentsline{toc}{chapter}{参考文献}

\begin{enumerate}
\def\labelenumi{(\Roman{enumi})}
\item
  J. von NEUMANN, Zur Operatorenmethode in der klassischen Mechanik, Ann. of Math., (2), 33 (1932), pp.~587--642.
\item
  A. KOLMOGOROFF, Grundbegriffe der Wahrscheinlichkeitsrechnung, Berlin (Springer), 1933.
\item
  I. GELFAND, Abstrakte Funktionen und lineare Operatoren, Mat. Sborn., (N.S.), 4 (1938), pp.~235--284.
\item
  N. DUNFORD, Uniformity in linear spaces, Trans. Amer. Math. Soc., 44 (1938), pp.~305--356.
\item
  N. DUNFORD AND B. PETTIS, Linear operations on summable functions, Trans. Amer. Math. Soc., 47 (1940), pp.~323--392.
\end{enumerate}

\chapter*{记号索引}
\phantomsection
\addcontentsline{toc}{chapter}{记号索引}

引用编号依次表示章、节、小节（在个别情形为习题）。

第 III 章 :

\(\mathcal{C}(\mathrm{X};\mathrm{E}),\mathcal{C}(\mathrm{X}),\mathcal{C}(\mathrm{X},\mathrm{A};\mathrm{E}),\mathcal{K}(\mathrm{X};\mathrm{E}),\mathcal{K}(\mathrm{X}),\mathcal{K}(\mathrm{X},\mathrm{A};\mathrm{E}),\mathcal{K}(\mathrm{X},\mathrm{A}),\mathcal{K}_{+}(\mathrm{X})\) （X 为局部紧空间，E 为拓扑向量空间）：III, 1, 1. Supp(f)（f 为取值于向量空间或 \(\overline{\mathbf{R}}\) 中的函数）：III, 1, 1。 \(\mathcal{C}^b (\mathrm{X};\mathrm{E}),\mathcal{C}^0 (\mathrm{X};\mathrm{E}):III,1,2.\) \(\| \mathbf{f}\|\) （f 为取值于赋范空间的函数）：III, 1, 2。 \(\mu (f),\langle f,\mu \rangle ,\int fd\mu ,\int f\mu ,\int f(x)d\mu (x),\int f(x)\mu (x)\) （f 为 \(\mathcal{H}(\mathrm{X};\mathbf{C})\) 中的函数， \(\mu\) 为（复）测度）：III, 1, 3。 \(\mathcal{M}(\mathrm{X};\mathbf{C}),\mathcal{M}(\mathrm{X}),\mathcal{M}_{\mathfrak{S}}(\mathrm{X};\mathbf{C}),\mathcal{M}_{\mathfrak{S}}(\mathrm{X}):III,1,3.\) \(\varepsilon_{a}:III,1,3.\) \(g\cdot \mu\) （g 为 \(\mathcal{C}(\mathrm{X};\mathbf{C})\) 中的函数）：III, 1, 4。 \(\overline{\mu},\mathcal{R}\mu ,\mathcal{I}\mu :III,1,5.\) \(\mathcal{M}(\mathrm{X};\mathbf{R}),\mathcal{M}(\mathrm{X}),\mathcal{M}_{+}(\mathrm{X}):III,1,5.\) \(\mu \leqslant \nu\) （ \(\mu ,\nu\) 为实测度）：III, 1, 5。 \(\mu^{+},\mu^{-},|\mu |\) （ \(\mu\) 为实测度）：III, 1, 5。 \(|\mu |\) （ \(\mu\) 为复测度）：III, 1, 6。 \(\| \mu \|\) （ \(\mu\) 为测度）：III, 1, 8。 \(\mathcal{M}^1 (\mathrm{X},\mathbf{R}),\mathcal{M}^1 (\mathrm{X}):III,1,8.\) \(\mu |Y\) （ \(\mu\) 为 X 上的测度，Y 为 X 的开子空间）：III, 2, 1。\\
Supp(μ)（ \(\mu\) 为测度）：III, 2, 2。 \(\langle \mathbf{f},\mathbf{z}'\rangle :III,3,1.\) \(\widetilde{\mathcal{K}} (\mathrm{X};\mathrm{E}):III,3,1.\) \(\int \mathbf{f}d\mu ,\int \mathbf{f}\mu ,\int \mathbf{f}(x)d\mu (x),\int \mathbf{f}(x)\mu (x)\) （ \(f\) 为 \(\widetilde{\mathcal{K}} (\mathrm{X};\mathrm{E}))\) 中的函数）：III, 3, 1。 \(\int d\mu (y)\int f(x,y)d\lambda (x):III,4,1.\) \(\iint f d\lambda d\mu ,\iint f d\mu d\lambda ,\iint f\lambda \mu ,\iint f\mu \lambda ,\iint f(x,y)d\lambda (x)d\mu (y),\) \(\iint f(x,y)d\mu (y)d\lambda (x),\iint f(x,y)\lambda (x)\mu (y),\iint f(x,y)\mu (y)\lambda (x):III,4,1.\)

\(\lambda \otimes \mu (\lambda ,\mu \text{为测度}):\mathrm{III},4,2.\) \(\mu_1\otimes \mu_2\otimes \dots \otimes \mu_n,\bigotimes_{i = 1}^n\mu_i:\mathrm{III},4,4.\) \(\int f d\mu_1d\mu_2\dots d\mu_n,\iint \dots \int f d\mu_1d\mu_2\dots d\mu_n,\int f(\mu_1\otimes \mu_2\otimes \dots \otimes \mu_n),\) \(\iint \dots \int f(x_1,x_2,\ldots ,x_n)d\mu_1(x_1)d\mu_2(x_2)\dots d\mu_n(x_n),\) \(\iint \dots \int f(x_1,x_2,\ldots ,x_n)\mu_1(x_1)\mu_2(x_2)\dots \mu_n(x_n):\mathrm{III},4,4.\) \(\bigotimes_{\lambda \in L}\mu_{\lambda}:\mathrm{III},4,6.\)

\(\varphi_{\mathrm{A}}:\mathrm{IV},1,1.\) \(\mathcal{K}_{+},\mathcal{I}_{+}(\mathrm{X}),\mathcal{I}_{+}:\mathrm{IV},1,1.\) \(\mu^{*}(f)\) （\(\mu\) 为正测度）：IV, 1, 1, IV, 1, 3 与 IV, 4, 习题 5。

\(\int^{*}fd\mu,\int^{*}f\mu,\int^{*}f(x)d\mu(x),\int^{*}f(x)\mu(x)\) （ \(f\) 为函数 \(\geqslant 0\) ， \(\mu\) 为正测度）：IV, 1, 3。

\(\mu^{*}(\mathrm{A})\) （A 为 X 的子集， \(\mu\) 为正测度）：IV, 1, 2 与 IV, 1, 4。

\(\widetilde{\mathbf{f}}:\mathrm{IV},2,4\) 与 IV, 2, 5.

\(\varphi((\widetilde{\mathbf{f}}_n))\), \(\widetilde{\mathbf{f}} +\widetilde{\mathbf{g}}\), \(\alpha \widetilde{\mathbf{f}},\widetilde{g}\widetilde{\mathbf{f}}:\mathrm{IV},2,4.\) \(\widetilde{f}\leqslant \widetilde{g},\sup_{n}\widetilde{f}_n,\inf_n\widetilde{f}_n,\limsup_{n\to \infty}\widetilde{f}_n,\liminf_{n\to \infty}\widetilde{f}_n\) （\(f,g,f_n\) 为数值函数）：IV, 2, 6。

\(|\mathbf{z}|\) （z 为赋范空间中的点）：IV, 3, 2。

\(|\mathbf{f}|\) （f 为取值于赋范空间的函数）：IV, 3, 2。

\(\mathrm{N}_p(\mathbf{f},\mu),\mathrm{N}_p(\mathbf{f}),\mathrm{N}_p(\widetilde{\mathbf{f}})\) （\(1\leqslant p<+\infty\)）：IV, 3, 2。

\(\mathcal{F}_{\mathrm{F}}(\mathrm{X}),\mathcal{F}_{\mathrm{F}}:\mathrm{IV},3,3.\) \(\mathcal{F}_{\mathrm{F}}^{p}(\mathrm{X},\mu),\mathcal{F}_{\mathrm{F}}^{p}(\mu),\mathcal{F}_{\mathrm{F}}^{p}:\mathrm{IV},3,3.\) \(\mathcal{N}_{\mathrm{F}}:\mathrm{IV},3,3.\) \(\mathcal{K}_{\mathrm{F}},\mathcal{L}_{\mathrm{F}}^{p}(\mathrm{X},\mu),\mathcal{L}_{\mathrm{F}}^{p}(\mu),\mathcal{L}_{\mathrm{F}}^{p},\mathrm{L}_{\mathrm{F}}^{p}(\mathrm{X},\mu),\mathrm{L}_{\mathrm{F}}^{p}(\mu),\mathrm{L}_{\mathrm{F}}^{p},\mathcal{L}^{p},\mathrm{L}^{p}\) （\(1\leqslant p<+\infty\)）：IV, 3, 4。

\(\| \widetilde{\mathbf{f}} \| _p\) （\(1\leqslant p < + \infty\)）：IV, 3, 4。

\(\mu (\mathbf{f}),\int \mathbf{f}d\mu ,\int \mathbf{f}(x)d\mu (x),\int \mathbf{f}\mu ,\int \mathbf{f}(x)\mu (x),\mu (\widetilde{\mathbf{f}})\) （f 为 \(\mu\) -可积的、取值于 Banach 空间的函数）：IV, 4, 1。

\(\mu (\mathrm{A})\) （A 为 \(\mu\) -可积集）：IV, 4, 5。

\(\mathcal{C}'(\mathrm{X};\mathbf{C}):\mathrm{IV},4,8.\) \(\mathcal{E}(\Phi),\mathcal{E}_{\mathrm{F}}(\Phi)\) （\(\Phi\) 为一个集环）：IV, 4, 9。

\(\mu_{*}(f)\) （\(f\) 为函数）：IV, 4, 习题 5。

\[
\mu_ {*} (\mathrm{A}) (\text { A a set }): \mathrm{IV}, 4, \text { Exer. } 7.
\]

\[
\int_ {\mathrm{A}} \textbf {f} d \mu , \int_ {\mathrm{A}} \textbf {f} \mu , \int_ {\mathrm{A}} ^ {*} f d \mu , \int_ {\mathrm{A}} ^ {*} f \mu : \mathrm{IV}, 5, 6.
\]

\[
\mathcal {S} (\mathrm{A}, \mu ; \mathrm{F}), \mathcal {S} _ {\mathrm{F}} (\mathrm{A}, \mu), \mathcal {S} _ {\mathrm{F}} (\mu), \mathcal {S} _ {\mathrm{F}}: \mathrm{IV}, 5, 11.
\]

\[
\boldsymbol {W} (\mathrm{V}, \mathrm{B}, \delta): \mathrm{IV}, 5, 11.
\]

\[
\mathrm{S} (\mathrm{A}, \mu ; \mathrm{F}), \mathrm{S} _ {\mathrm{F}} (\mathrm{A}, \mu), \mathrm{S} _ {\mathrm{F}} (\mu), \mathrm{S} _ {\mathrm{F}}: \mathrm{IV}, 5, 11.
\]

\[
\mathrm{M} _ {\infty} (f), m _ {\infty} (f): \mathrm{IV}, 6, 2.
\]

\[
\mathrm{N} _ {\infty} (\mathbf {f}): \mathrm{IV}, 6, 3.
\]

\[
\mathcal {L} _ {\mathrm{F}} ^ {\infty} (\mathrm{X}, \mu), \mathcal {L} _ {\mathrm{F}} ^ {\infty} (\mu), \mathcal {L} _ {\mathrm{F}} ^ {\infty}, \mathcal {N} _ {\mathrm{F}} ^ {\infty}: \mathrm{IV}, 6, 3.
\]

\[
\| \dot {\mathbf {f}} \| _ {\infty}, \mathrm{N} _ {\infty} (\dot {\mathbf {f}}): \mathrm{IV}, 6, 3.
\]

\[
\mathrm{L} _ {\mathrm{F}} ^ {\infty} (\mathrm{X}, \mu), \mathrm{L} _ {\mathrm{F}} ^ {\infty} (\mu), \mathrm{L} _ {\mathrm{F}} ^ {\infty}, \mathcal {L} ^ {\infty}, \mathrm{L} ^ {\infty}: \mathrm{IV}, 6, 3.
\]

\[
\mathbf {b} _ {\mu}: \mathrm{IV}, 7, 1.
\]

\[
\mathrm{Ch} _ {\mathcal {H}} (\mathrm{X}), \mathrm{Ch} (\mathrm{X}), \check {\mathrm{S}} _ {\mathcal {H}} (\mathrm{X}), \check {\mathrm{S}} (\mathrm{X}): \mathrm{IV}, 7, 3.
\]

\[
\mu^ {\bullet} (f), \mu^ {\bullet} (\mathrm{A}), \int^ {\bullet} f d \mu , \int^ {\bullet} f (t) d \mu (t), \int^ {\bullet} f \mu : \mathrm{V}, 1, 1.
\]

\[
\overline {{\mathcal {F}}} _ {\mathrm{F}} ^ {p} (\mathrm{T}, \mu), \overline {{\mathcal {F}}} _ {\mathrm{F}} ^ {p} (\mu), \overline {{\mathcal {F}}} _ {\mathrm{F}} ^ {p}: \mathrm{V}, 1, 3.
\]

\[
\overline {{\mathbf {N}}} _ {p} (f), \overline {{\mathcal {L}}} _ {\mathrm{F}} ^ {p} (\mathrm{T}, \mu), \overline {{\mathcal {L}}} _ {\mathrm{F}} ^ {p} (\mu), \overline {{\mathcal {L}}} _ {\mathrm{F}} ^ {p}: \mathrm{V}, 1, 3.
\]

\[
\overline {{\mathcal {L}}} _ {\mathrm{F}} ^ {p} (\mathrm{T}, \theta) (\theta \text {a complex measure}): \mathrm{V,1,3.}
\]

\[
\int \lambda_ {t} d \mu (t) (t \mapsto \lambda_ {t} \text {a family of positive measures}): \mathrm{V}, 3, 1.
\]

\[
\int d \mu (t) \int f (x) d \lambda_ {t} (x): \mathrm{V}, 3, 1.
\]

\[
\| \Lambda \| (\Lambda \text {a diffusion}): \mathrm{V}, 3, 5.
\]

\[
\langle \eta , h \rangle : \mathrm{V}, 3, 5.
\]

\[
\Lambda f, \mu \Lambda : \mathrm{V}, 3, 5.
\]

\[
\Lambda \mathrm{H}: \mathrm{V}, 3, 6.
\]

\[
\mathcal {L} _ {\mathrm{loc}} ^ {1} (\mathrm{T}, \mu ; \mathrm{F}), \mathrm{L} _ {\mathrm{loc}} ^ {1} (\mathrm{T}, \mu ; \mathrm{F}): \mathrm{V}, 5, 1.
\]

\[
\int_ {\mathrm{A}} ^ {\bullet} f d \mu : \mathrm{V}, 5, 3.
\]

\[
u (\mu_ {1}, \dots , \mu_ {n}) (u \text {a positively homogeneous numerical function}): \mathrm{V}, 5, 9.
\]

\[
\pi^ {- 1} (\mu) (\pi \text {a local homeomorphism}) \colon \mathrm{V,6,6.}
\]

\[
\iint^ {*} f (t, t ^ {\prime}) d \mu (t) d \mu^ {\prime} (t ^ {\prime}), \iint^ {\bullet} f (t, t ^ {\prime}) d \mu (t) d \mu^ {\prime} (t ^ {\prime}), \iint \mathbf {f} (t, t ^ {\prime}) d \mu (t) d \mu^ {\prime} (t ^ {\prime}): \mathrm{V}, 8, 1
\]

第 VI 章 :

\(\mathbf{F}^{\prime},\mathbf{F}^{\prime \prime},\mathbf{F}^{\prime *},\mathbf{F}_{\sigma}\) （F 为 Hausdorff 局部凸空间）：VI, 引言。

\(\mathcal{K}(\mathrm{T}),\mathcal{K}_{\mathbf{R}}(\mathrm{T}),\mathcal{K}_{\mathbf{C}}(\mathrm{T}),\mathcal{K}(\mathrm{T},\mathrm{A}),\mathcal{K}_{\mathbf{C}}(\mathrm{T},\mathrm{A}):\mathrm{VI},\) 引言。

\(\langle \mathbf{f},\mathbf{z}^{\prime}\rangle ,\langle \mathbf{z}^{\prime},\mathbf{f}\rangle :\mathrm{VI},1.\)

\(\int \mathbf{f} d\mu, \int \mathbf{f}(t) d\mu(t)\) （f 为向量值函数， \(\mu\) 为正测度）：VI, 1, 1。

\(g\mathbf{f},\mathbf{f}g\) （f 为向量值函数， \(g\) 为标量函数）：VI, 1, 1。

\(\mathcal{C}'(\mathrm{T}) : \mathrm{VI}, 1, 6.\)

\(\int f d\mathbf{m}, \int f(t) d\mathbf{m}(t)\) （ \(f\) 为数值函数， \(\mathbf{m}\) 为向量测度）：VI, 2, 1 与 VI, 2, 2。

\(g \cdot m\) （g 为数值函数，m 为向量测度）：VI, 2, 1。

\(\mathcal{L}(\mathbf{m}) : \mathrm{VI}, 2, 2.\)

\(q(\mathbf{m}), |\mathbf{m}|\) （ \(q\) 为半范数， \(\mathbf{m}\) 为向量测度）：VI, 2, 3。

\(\mathbf{f} \cdot \mu\) （f 为向量值函数， \(\mu\) 为正测度）：VI, 2, 4。

\(\mathcal{L}_{\mathrm{F}_s'}^{\infty}, \mathrm{L}_{\mathrm{F}_s'}^{\infty} : \mathrm{VI}, 2, 5.\)

\(\langle \mathbf{f},\mathbf{g}\rangle\) （f,g 为向量值函数）：VI, 2, 6。

\(\mathrm{I}_{\Phi ,\mathbf{m}},\int \mathbf{f}dm\) （f 为向量值函数，m 为向量测度）：VI, 2, 7。

\(|m|, \int \mathbf{f} dm \ (m \ a \ complex \ measure) : VI, 2, 8, III, 1, 6.\)

\(\mathcal{L}_{\mathrm{F}}^{p}(\mathrm{T},m), \overline{\mathcal{L}}_{\mathrm{F}}^{p}(\mathrm{T},m), \mathrm{L}_{\mathrm{F}}^{p}(\mathrm{T},m) (m \text{ 为复测度}): \mathrm{VI}, 2, 8, \mathrm{V}, 1, 3.\)

\(h \cdot m\) （ \(m\) 为复测度）：VI, 2, 8。

\(\overline{m}\) （ \(m\) 为复测度）：VI, 2, 8, III, 1, 5。

\(\| m\|\) （ \(m\) 为复测度）：VI, 2, 9, III, 1, 8。

\(\pi (m), m_{\mathrm{Y}}, m \otimes m' (m, m'\) 为复测度）：VI, 2, 10。

\(\mathfrak{B}(\mathrm{F}_{1},\mathrm{F}_{2}),{}^{r}\Phi,{}^{l}\Phi :\mathrm{VI},\mathrm{App.},1.\)

\(\mathrm{E}_{\sigma}, \mathrm{F}_{\sigma}, \mathrm{E}_{s}^{\prime}, \mathrm{F}_{s}^{\prime}, \mathcal{B}(\mathrm{E}, \mathrm{F}) : \mathrm{VI}, \mathrm{App.}, 1.\)

\(\Lambda_{\mathrm{F}^{\prime}}^{p}(\mathrm{T},\mu),\mathrm{M}_{p},\mathrm{M}_{p}^{\prime}:\mathrm{VI},1,\mathrm{Exer.}16.\)
